\documentclass[
  11pt,					% Schriftgröße
  DIV=13,				% Seitenlayout (Satzspiegel)
  parskip=half,			% Abstand zwischen Absätzen
  twoside,				% Doppelseitig
%  headsepline,
  headings,	
%  draft,				% Korrek\add{text}turfassung
  ]{scrreprt}			% scrartcl	


% Kotlin highlighting
\usepackage[dvipsnames]{xcolor}
\usepackage{listings}
\lstdefinelanguage{Kotlin}{
	comment=[l]{//},
	commentstyle={\color{darkgray}\ttfamily},
	emph={filter, first, firstOrNull, forEach, lazy, map, mapNotNull, println},
	emphstyle={\color{OrangeRed}},
	identifierstyle=\color{black},
	keywords={!in, !is, abstract, actual, annotation, as, as?, break, by, catch, class, companion, const, constructor, continue, crossinline, data, delegate, do, dynamic, else, enum, expect, external, false, field, file, final, finally, for, fun, get, if, import, in, infix, init, inline, inner, interface, internal, is, lateinit, noinline, null, object, open, operator, out, override, package, param, private, property, protected, public, receiveris, reified, return, return@, sealed, set, setparam, super, suspend, tailrec, this, throw, true, try, typealias, typeof, val, var, vararg, when, where, while, returns, variable, function},
	keywordstyle={\color{NavyBlue}\bfseries},
	morecomment=[s]{/*}{*/},
	morestring=[b]",
	morestring=[s]{"""*}{*"""},
	ndkeywords={@Deprecated, @JvmField, @JvmName, @JvmOverloads, @JvmStatic, @JvmSynthetic, Array, Byte, Double, Float, Int, Integer, Iterable, Long, Runnable, Short, String, Any, Unit, Nothing},
	ndkeywordstyle={\color{BurntOrange}\bfseries},
	sensitive=true,
	stringstyle={\color{ForestGreen}\ttfamily}
}

\lstset{
	basicstyle=\scriptsize\sffamily\color{black},
	frame=single,
	numbers=left,
	numbersep=5pt,
	numberstyle=\tiny\color{gray},
	showspaces=false,
	showstringspaces=false,
	tabsize=4,
	linewidth=\textwidth,
	captionpos=b,
	extendedchars=true
	inputencoding=utf8,
	literate={ö}{{\"o}}1
	{ä}{{\"a}}1
	{ü}{{\"u}}1
	{Ö}{{\"O}}1
	{Ä}{{\"A}}1
	{Ü}{{\"U}}1
}

\usepackage[utf8]{inputenc}
\usepackage[ngerman,english]{babel}
\selectlanguage{ngerman}

\usepackage[ddmmyyyy]{datetime}
\renewcommand{\dateseparator}{.}

\usepackage{graphicx}
\graphicspath{{bilder/}}
\usepackage{svg}
\usepackage{float} % Force placement of figures with [H]

\usepackage{enumitem}
\usepackage{tikz}

% Blindtext
\usepackage{blindtext} %TODO: Remove later
\let\oldblindtext\blindtext
\renewcommand{\blindtext}{{\color{red}\oldblindtext}}
 
% Schrifteinstellungen
\usepackage{lmodern} 		% Vektorschrift
\renewcommand{\familydefault}{\sfdefault}
\usepackage{sansmath}
\sansmath 
\usepackage{microtype}

% Literatur einbinden
\usepackage{csquotes}	% Steuerung der Anführungszeichen
\usepackage[
  backend=bibtex,		% Sortier-Compiler
  style=numeric,	% Zitationsstil
  sorting=none
  ]{biblatex}
  
\addbibresource{ref/ref.bib} %Citavi Reference Datei

% Mathemodus
\usepackage{amsmath,amssymb}

% Ensure figures are placed correctly
\usepackage[section]{placeins}

% Trennung
\hyphenation{Crash-zo-ne}



% Kopf- und Fußzeile
\usepackage[
headsepline,		% Kopfzeilen-Sepparationslinie
automark,		% Lebende Kolumnentitel
]
{scrlayer-scrpage}
\pagestyle{scrheadings}		
\ofoot*{\pagemark}			
\ohead*{\headmark}

\usepackage{abstract}
\usepackage{adjustbox}


\usepackage[acronym]{glossaries-extra}
\setabbreviationstyle[abbreviation]{short-desc}
\makeglossaries

\def\table{\def\figurename{Tabelle}\figure}
\let\endtable\endfigure

\usepackage{caption}
\usepackage{newfloat}

%%%%%%%%%%%%%%%%%%%%%%%%%%%%%%%%%%%%%%%%%%%%%%%%%
%%%%%%%%%%%%%%%%%%%%%%%%%%%%%%%%%%%%%%%%%%%%%%%%%
%%%%%%%%%%%%%%%%%%%%%%%%%%%%%%%%%%%%%%%%%%%%%%%%%
%%%%%%%%%%%%%%%%%%%%%%%%%%%%%%%%%%%%%%%%%%%%%%%%%
%%%%%%%%%%%%%%%%%%%%%%%%%%%%%%%%%%%%%%%%%%%%%%%%%
%%%%%%%%%%%%%%%%%%%%%%%%%%%%%%%%%%%%%%%%%%%%%%%%%
%%%%%%%%%%%%%%%%%%%%%%%%%%%%%%%%%%%%%%%%%%%%%%%%%
%%%%%%%%%%%%%%%%%%%%%%%%%%%%%%%%%%%%%%%%%%%%%%%%%

% Titelseite
\titlehead{
  \begin{center}
  	\includegraphics[width=0.5\textwidth]{thi_logo_cropped}
  \end{center}
}

\title{Aggregationsschnitt einer Checkout-Software auf Basis einer Hexagonalen Architektur mit Domain-Driven Design }

\subtitle{ \vspace{2ex} \LARGE Bachelor-Arbeit}

\author{Simon Thalmaier}

\date{}

\publishers{
  \begin{tabular}{rl}
   \textbf{Erstprüfer} 		& Prof. Dr. Sebastian Apel	\\
   \textbf{Zweitprüfer} 	& Prof. Dr. rer. nat. Franz Regensburger	\\
   \textbf{Betreuer} 		& Stefano Lucka  	\\
   \textbf{Ausgabedatum} 	& - 				\\
   \textbf{Abgabedatum} 	& -					\\
  \end{tabular}
  }
  
% Rückseite der Titelseite
\uppertitleback{
	
\adjustbox{minipage=0.8\textwidth}{
\vspace{1.5cm} 
	{\Large\textbf{Erklärung zur Bachelorarbeit}}
	
	\vspace{2em}
	
	{
		
		Ich erkläre hiermit, dass ich diese Bachelorarbeit selbständig verfasst,
		noch nicht anderweitig für Prüfungszwecke vorgelegt, keine anderen als
		die angegebenen Quellen und Hilfsmittel benutzt sowie wörtliche und
		sinngemäße Zitate als solche gekennzeichnet habe. 
		
		\vspace{5em}
		
		\begin{tikzpicture}[remember picture,overlay]
			\node[xshift=87mm,yshift=-69mm, anchor=north west] at (current page.north west){%
				\includegraphics[width=0.5\textwidth]{Unterschrift.png}};
		\end{tikzpicture}
		
		
		
		Ingolstadt, \today
		\hfill
		\rule{220pt}{1pt} \\
		
		\hfill
		Simon Thalmaier 
	}
\par}
}

\lowertitleback{} %Dokumenteninformation, Veröffentlichung, Rahmen, bibliographisches Angaben



\newcommand\frontmatter{%
	\cleardoublepage
	%\@mainmatterfalse
	\pagenumbering{Roman}}

\newcommand\mainmatter{%
	\cleardoublepage
	% \@mainmattertrue
	\pagenumbering{arabic}}

\newcommand\backmatter{%
	\cleardoublepage
	%\@mainmatterfalse
	\pagenumbering{roman}}

\newcommand{\comment}[1]{{\color{Green}//Kommentar: #1}\\} %TODO: Delete later
\newcommand{\note}[2]{{\color{red}#1}}
\newcommand{\domainModell}[2]{
	{\large \textbf{#1:}}
	\begin{itemize}[topsep=-5pt]#2\end{itemize}
	\vspace{12pt}
}
\newcommand{\domainModellWithoutTitle}[1]{
	\begin{itemize}[topsep=-5pt]#1\end{itemize}
	\vspace{12pt}
}
\newcommand{\subDomainModell}[1]{
	\begin{itemize}[noitemsep,topsep=-5pt]#1\end{itemize}
	\vspace{8pt}
}

\usepackage{lscape}
\newcommand{\centertitle}[1]{\vspace{0.8mm}#1\vspace{0.5mm}}
\newcommand{\centertable}[1]{\vspace{0.1cm}#1\vspace{0.1cm}}

\usepackage{array} 

\newcommand{\ul}[1]{\emph{#1}}

\begin{document}
  
  
\newacronym{POC}{POC}{Proof-of-Concept}
\newacronym{DDD}{DDD}{Domain-Driven Design}
\newacronym{SRP}{SRP}{Single-Responsibility-Prinzip}
\newacronym{OCP}{OCP}{Open-Closed-Prinzip}
\newacronym{LSP}{LSP}{Liskovsches Substitutionsprinzip}
\newacronym{ISP}{ISP}{Interface-Segregation-Prinzip}
\newacronym{DIP}{DIP}{Dependency-Inversion-Prinzip}
\newacronym{DTO}{DTO}{Data-Transfer-Object}
\newacronym{HTTP}{HTTP}{Hypertext Transfer Protocol}
\newacronym{REST}{REST}{Representational State Transfer}
\newacronym{CRUD}{CRUD}{Create Update Delete}
\newacronym{CQRS}{CQRS}{Command and Query Responsibility Segregation}
\newacronym{KPI}{KPI}{Key-Performance-Indicator}

\newglossaryentry{Serialisierung}{name={Serialisierung}, %TODO: Quelle? https://www.duden.de/rechtschreibung/serialisieren
	description={Konvertierung von Datenobjekte in ein sequenzielles Format}
}
\newglossaryentry{Deserialisierung}{name={Deserialisierung}, 
	description={Die Umkehrung einer Serialisierung. Wandelt ein sequenzielles Datenformat in Objekte einer Klasse um.}
}
\newglossaryentry{User Story}{name={User Story},
	description={Ein möglicher Ablauf des kompletten Businessprozesses aus Sicht des Kunden.}
}
\newglossaryentry{Race Condition}{name={Race Condition},
	description={Eine Situtation in welcher, aufgrund von parallelen Aktionen, das Resultat von deren zeitlichen Stattfinden abhängig ist.}}
\newglossaryentry{Boilerplate}{name=Boilerplate,
	description={Ein Teil einer Software, welcher viele Zeilen an Code einnimmt, obwohl dadurch nur wenig bis gar keine Funktion bereitgestellt wird}
}
\newglossaryentry{Stakeholder}{name=Stakeholder,
	description={Eine Gruppe von Personen mit relevanten Interesse und Einfluss auf eine Sache bzw. Projekt}
}
\newglossaryentry{Lost Update}{name=Lost Update,
	description={Phänomen, welches bei zeitgleichen Operationen auf den gleichen Datensätzen auftreten kann. Die angepassten Datensätze einer Transaktion gehen verloren, da sie direkt von einer zweiten Transaktion überschrieben werden, welche jedoch als Ausgangspunkt noch auf den alten Stand durchgeführt worden ist}
}
\newglossaryentry{Information-Expert-Prinzip}{name=Information-Expert-Prinzip,
	description={Die Verantwortung eines Anliegens liegt bei der Komponente, welche die notwendigen Informationen zur Erfüllen besitzt}
}
\newglossaryentry{Invariante}{name=Invariante,
	description={Bedinung, welche auch nach Datenanpassungen jederzeit erfüllt sein muss}
}
\newglossaryentry{Kohasion}{name=Kohäsion,
	description={Grad der logischen inneren Zusammengehörigkeit einer Komponente. Komponente, welche nur eng beinaheliegende Aufgaben erfüllen, haben einen hohen Grad an Kohäsion}
}
\newglossaryentry{DI}{name={Dependency Injection},
	description={Eine erweiterte Form des Dependency-Inversion-Prinzip, welches Abhängigkeiten erst zur Laufzeit des Programmes hinzufügt}
}
\newglossaryentry{immutable}{name=immutable,
	description={Die Unveränderlichkeit von Werten bzw. Variablen innerhalb einer Klasse}
}
\newglossaryentry{Scrum}{name=Scrum,
	description={Eine agile Entwicklungsmethode, welches hohen Fokus auf kontinuierliche Verbesserung in einem geregelten Zyklus legt}
}
\newglossaryentry{Sprint}{name=Sprint,
	description={Ein wiederkehrender festgelegter Zeitraum, indem eine vorher definierter Umfang an Arbeitspakten abgearbeitet wird}
}

\newglossaryentry{Product Owner}{name=Product Owner,
	description={Eine Scrum-Rolle, welche die wirtschaftliche Ziele und Prioritäten der Aufgabenpakete bestimmt}
}
\newglossaryentry{Lazy Loading}{name={Lazy Loading},
	description={Das Laden von Daten aus einem Datenspeicher oder sonstigen Quellen wird erst durchgeführt, sobald diese benötigt werden, wodurch unnötiges Laden minimiert wird.}}
\newglossaryentry{Collection}{name={Collection},
	description={Analog zu den Tabellen in einer relationalen Datenbank, stellen persistieren Collections pro Eintrag einen Datensatz.}}
\newglossaryentry{Connection-Pool}{name={Connection-Pool},
	description={Eine Ansammlung von Verbindungen meist zu Datenbanken zur Performance-Optimierung, indem alte Verbindungen wiederverwendet werden.}
}
%\newglossaryentry{Enumeration}{name=Enumeration,
%	description={Eine Auflistung von konstanten, unveränderlichen Werten}
%}
    
  \maketitle
  
  \frontmatter
  
  
%% ABSTRACT %%


{
	
	\raggedbottom
	\centering
	
	\selectlanguage{ngerman}
	
	\large
	{
		\vspace*{\fill}
		\adjustbox{minipage=0.88\textwidth}{
			\begin{center}
				\textbf{Erklärung zur Bachelorarbeit}
			\end{center}
			\vspace{1em}  
			
			Ich erkläre hiermit, dass ich diese Bachelorarbeit selbständig verfasst,
			noch nicht anderweitig für Prüfungszwecke vorgelegt, keine anderen als
			die angegebenen Quellen und Hilfsmittel benutzt sowie wörtliche und
			sinngemäße Zitate als solche gekennzeichnet habe. 
			
			\vspace{4em}
			
			\begin{tikzpicture}[remember picture,overlay]
				\node[xshift=95mm,yshift=-136mm, anchor=north west] at (current page.north west){%
					\includegraphics[width=0.5\textwidth]{Unterschrift.png}};
			\end{tikzpicture}
			
			Ingolstadt, \today
			\hfill
			\rule{220pt}{1pt} \\
			
			\hfill
			Simon Thalmaier 
			
			\par}
		\vfill
	}
	
	
	
	\pagebreak
	
	
	
	\selectlanguage{english}
	\large
	
	{
		\vspace*{\fill}
		\adjustbox{minipage=0.88\textwidth}{
			\begin{center}
				\textbf{Abstract}
			\end{center}
			\vspace{1em}  
			
			In the e-commerce sector software projects fulfill complex business requirements and therefore need a stable architecture. This thesis examines an online shop checkout domain and how the bounded context can be implemented utilizing hexagonal architecture and domain-driven design. The focus is placed on the design of the aggregates. Various data models are analyzed and evaluated based on their complexity, performance, parallelism and client-friendliness. Generally, big aggregates suffer from lowered parallelism, however can be implemented more easily since all required information are loaded at the same time from the database. Splitting the data model into different aggregates entails the use of eventual consistency or transaction over more than one aggregate. Eventual consistency rises the complexity of the checkout software, on the other hand cross-aggregate transactions make the usage of multiple database hosts more difficult. The performance of the individual aggregate designs is measured with a load test. On average applications with one cohesive aggregate process more requests per second than a model using separated aggregates. Conclusively, the checkout software profits from a higher performance and reduced complexity by implementing only one aggregate. If a common use case is the mutation of one resource by distinct processes simultaneously then the affected objects need to be placed in different aggregates. 
			
			\par}
		\vfill
	}
	
	
 	\pagebreak

	\selectlanguage{ngerman}
	\large
	
	{
		\vspace*{\fill}
		\adjustbox{minipage=0.88\textwidth}{
			\begin{center}
				\textbf{Zusammenfassung}
			\end{center}
			
			\vspace{1em}  
			
			Softwareprojekte im E-Commerce-Bereich erfüllen komplexe Businessanforderungen und benötigen aufgrund dessen eine stabile Architektur. Dieses Projekt untersucht die Domain eines Onlineshop-Checkouts und wie der Bounded-Context mithilfe einer Hexagonalen Architektur und Domain-Driven Design implementiert werden kann. Besonders liegt der Aggregationsschnitt im Fokus, wobei unterschiedliche Datenmodelle analysiert und anhand von Komplexität, Performance, Parallelität und Client-Freundlichkeit bewertet werden. Große Aggregates leiden generell unter verringerter Parallelität, jedoch bietet ein zusammengehöriges Datenmodell eine vereinfachte Umsetzung von Businessanforderungen, da stets alle Informationen aus der Datenbank geladen werden. Die Aufteilung in mehrere Aggregates erzwingt die Anwendung von Eventueller Konsistenz oder einer Transaktion über mehrere Aggregates. Eventuelle Konsistenz erhöht die Komplexität der Checkout-Software, wohingegen eine aggregate-übergreifende Transaktion die Verwendung von unterschiedlichen Datenbank-Hosts erschwert. Anhand eines Lasttests wird die Performance der Designansätze betrachtet. Im Durchschnitt verarbeitet die Applikation mit einem zusammengehörigen Datenmodell mehr Anfragen pro Sekunde als unter Verwendung getrennter Aggregates. Als Fazit dieser Arbeit wird argumentativ begründet, dass innerhalb einer Checkout-Software die Vorteile eines Designs mit einem einzigen Aggregate dank der erhöhten Performance und reduzierten Komplexität überwiegen. Falls eine zeitgleiche Bearbeitung einer Ressource durch unterschiedliche Prozesse ein gängiger Anwendungsfall ist, müssen die betroffenen Objekte in separate Aggregates verlagert werden.
			
			\par}
		\vfill
	}
	
	
	
	\pagebreak

	\selectlanguage{ngerman}
	\large
	
	{
		\vspace*{\fill}
		\adjustbox{minipage=0.88\textwidth}{
			\begin{center}
				\textbf{Danksagung}
			\end{center}
			\vspace{1em} 
			
			
			Mein Dank gilt der Firma MediaMarktSaturn Technology, welche mir ermöglicht hat als Werkstudent über die vergangenen drei Jahre zu arbeiten und das Thema für diese Bachelorarbeit bereitgestellt hat. Besonders bedanke ich mich bei Sebastian Jurjanz für die Unterstützung in meiner Ausbildung und während dieses Projektes. Zusätzlich ist die Bearbeitung des Forschungsthemas in diesem Umfange durch Stefano Lucka dank seiner Betreuung und seinem Feedback ermöglicht worden. \\
			
			Aufseiten der Technischen Hochschule Ingolstadt bedanke ich mich bei \text{Professor Dr. Sebastian Apel} für das konstruktive Feedback und die umfangreiche Beratung. 
			
			
			\par}
		\vfill
	}
	
		
	\pagebreak
	
\par}
  
  \selectlanguage{ngerman}
  
  \tableofcontents
  
  %Abbildungsverzeichnis
  \newpage
  \addcontentsline{toc}{chapter}{Darstellungsverzeichnis}
  \renewcommand{\listfigurename}{Darstellungsverzeichnis}
  \listoffigures
  
  %Abbildungsverzeichnis
  \newpage
  \renewcommand{\lstlistingname}{Codebeispiel}
  \renewcommand*{\lstlistlistingname}{Codebeispiel-Verzeichnis}
  \addcontentsline{toc}{chapter}{Codebeispiel-Verzeichnis}
  \lstlistoflistings
  
  %Acronyms
  \newpage
  \setlength{\glsdescwidth}{0.8\textwidth}
  \printglossary[type=\acronymtype, style=super, nonumberlist]
  
  %Glossary
  \newpage
  \setlength{\glsdescwidth}{0.75\textwidth}
  \printglossary[nonumberlist, style=super] %TODO: Disable sorting
  
  \mainmatter
  
  
%% EINLEITUNG %% 

\chapter{Einleitung}

Heutzutage werden Applikationen für den langfristigen Gebrauch in produktiver Umgebung entwickelt. Im Durchschnitt kann hierbei der erwartete Lebenszyklus dieser Anwendungen auf circa 10 Jahre festgelegt werden \cite{Tamai.1992}, weshalb sich über die vergangenen Jahrzehnte ein starker Fokus auf ein wartbares und flexibles Softwaredesign gebildet hat. Darüber hinaus erfüllen sie umfangreiche Anwendungsprofile und müssen diese fehlerfrei und performant bewältigen, wodurch ebenfalls ein stabiler Architekturansatz von Nöten ist \cite{Bosch.2001}. Dementsprechend stehen Software Engineers mittlerweile eine große Bandbreite an Entwurfsmustern und Anti-Pattern zur Verfügung, wie beispielsweise die Hexagonale Architektur, Event-Storming, \Gls{Anemic Domain Model} und Microservices. Letzteres gewann aufgrund ihrer Skalierbarkeit und losen Kopplung der Komponenten in den vergangenen Jahren zunehmend an Bedeutung \cite{oreilly.Microservices, Sampaio.2017}. Hiermit einhergehend, erfuhr das im Jahre 2011 erschiene Buch \citetitle{Evans.2011} von \citeauthor{Evans.2011} an Beliebtheit, da viele der bearbeiteten Themenschwerpunkte auf Microservices adaptiert werden können \cite[S. 130ff.]{Vernon.2015}\cite{Microservice.DDD.2017}. Domain-Driven Design stellt zum Lösen gängiger Problemstellungen der Softwareentwicklung unter anderem ein Vorgehensmodell für die Realisierung eines businessorientierten Datenmodells bereit, das bei der Umsetzung von Businessanforderungen unterstützt. Eine Checkout-Software, welche in dieser Bachelorarbeit als Proof-of-Concept entwickelt wird, profitiert von den Vorteilen einer solchen Architektur. \comment{Letzte zwei Sätze auch Quellen?}

\section{Problemstellung}

Ein elementarer Bestandteil der Funktionsweisen eines Onlineshops ist der Warenkorb. In diesem können zum späteren Erwerb Waren abgelegt oder zusätzliche Dienstleistungen hinzugefügt werden. Im Verlauf des Kaufprozesses ist es zudem möglich, Kundendaten zu hinterlegen, sowie die gewünschte Versandart und Zahlungsarten auszuwählen. Nach erfolgreicher Überprüfung mittels Validierungsrichtlinien findet die Kaufabwicklung statt, der sogenannte 'Checkout'. Um die zuvor genannten Anwendungsfälle zu verwirklichen, wird eine eigens dafür geschriebene Software benötigt. In dieser Bachelorarbeit wird diese Anwendung vereinfacht implementiert und als 'Checkout-Software' bezeichnet. Sie erfährt stetige Änderungen, beinhaltet im Vergleich zu anderen Softwareprojekten umfangreiche Businesslogik und ihre Antwortzeiten haben durch die Einbindung in das Frontend auch direkten Einfluss auf das Kundenerlebnis. Dadurch liegt stets ein Fokus auf die Erfüllung von Qualitätsmerkmalen, wie Stabilität, Testbarkeit und Wartbarkeit. Der Checkout-Prozess, welcher durch diese Software abgewickelt wird, muss für alle Länder, in welche die Applikation operiert, und ihre individuellen gesetzlichen Anforderungen fiskalisch korrekt ausgeführt werden. Jederzeit können neue Businessanforderungen entstehen, wodurch weitere länderspezifische Richtlinien in den Zuständigkeitsbereich der Anwendung fallen und beispielsweise eine Anpassung des Datenmodells erfordern. Zudem erfordert das System zur Abwicklung ihrer Arbeitsprozesse Daten aus verschiedensten Unternehmensbereichen wie Preise, Lieferkosten, Produkt- und Bestandsinformationen. Die Kommunikation mit externen Komponenten erhöht die Komplexität der Anwendung und erfordert die Berücksichtigung systemübergreifender Anforderungen. Eine große Rolle spielt hierbei die Performance. Vor allem bei hoher Auslastung, beispielsweise während Kampagnen, muss das Gesamtsystem weiterhin zuverlässig alle Anfrage in akzeptabler Zeit abarbeiten können. Dementsprechend stellt die Implementierung einer solchen Checkout-Applikation für Software Engineers eine große Herausforderung dar. Sofern die verwendete Architektur im langjährigen Entwicklungsprozess an Übersichtlichkeit verliert, leidet zugleich auch die Wartbarkeit des Sourcecodes. Als Folge können weitere Qualitätsmerkmale negativ betroffen sein und der Aufwand zur Umsetzung von Businessanforderungen steigt \cite[S. 3f.]{Pigoski.2001}. Daraus ergibt sich für den Checkout der grundsätzlicher Bedarf zur Befolgung bestimmter Industriestandards hinsichtlich der Softwarearchitektur und des Datenmodells. Somit besteht ein Teilaspekt der Problemstellung in der Auswahl einer geeigneten Architektur zur Realisierung der Software. 

Auf Basis der kommenden Kapitel wird die Verwendung einer Hexagonaler Architektur inklusive Domain-Driven Design für den Proof-of-Concept argumentativ begründet. Die Projektdurchführung orientiert sich hierbei am empfohlenen Entwicklungsprozess von Domain-Driven Design. Im Fokus der Bachelorarbeit stehen die möglichen Einteilungen des Datenmodells in Aggregates, in dieser Arbeit als 'Aggregationsschnitt' bezeichnet, welche durch die Untersuchung der Anwendungsfälle erschlossen werden. Das Forschungsthema bildet sich aus der Frage, welche funktionalen und nicht-funktionalen Implikationen die unterschiedliche Aggregationsschnitte auf die Applikationen besitzen. Hierzu werden sie in Form eines Proof-of-Concepts implementiert, analysiert und anhand von Performance-Tests bewertet.

\section{Das Unternehmen MediaMarktSaturn}

Dieses Projekt wurde in Zusammenarbeit mit dem Unternehmen \emph{MediaMarktSaturn Retail Group GmbH}, kurz \emph{MediaMarktSaturn}, erarbeitet. Mit über 1.000 Märkten sowie den Onlineshops in 13 Ländern ist MediaMarktSaturn Europas größter Anbieter von Unterhaltungselektronik sowie zugehöriger Dienstleistungen und Services. Die umfangreiche Produktauswahl in Kombination mit passenden Services und Kundennähe schaffen ein einzigartiges Einkaufserlaubnis - über alle Verkaufskanäle hinweg. Die Zugehörigkeiten der Märkte ist hierbei in die Marken \emph{Media Markt} und \emph{Saturn} unterteilt. \cite{mms.Unternehmen}

Wegen des massiv steigenden Onlineanteils gewann der Onlineshop über die Jahre für Media Markt und Saturn zunehmend an Bedeutung. Der Ausbruch der Corona-Pandemie und die damit verbundene europaweite Schließung der Märkte hat die Verlagerung der Umsatzeinnahmen vom Markt- zum Onlinegeschäft nochmals verschärft. Als Folge dessen wurden die Unternehmensziele auf die Entwicklung komplexer Software zur Unterstützung des Onlineshops neu ausgelegt. Die Umsetzung obliegt der zentralen IT-Gesellschaft MediaMarktSaturn Technology, wo über 700 interne und 1000 externe Engineers in einer skalierten Produktorganisation an der steten Optimierung der Systemlandschaft arbeiten \cite{mms.technology}.

Die Durchführung des Projektes bzw. Proof-of-Concepts geschah in Kooperation mit dem Bereich \emph{Checkout \& Payment}. Die sechs zugehörigen Teammitglieder sind zuständig einen unternehmensweiten, universellen Checkout für alle Länder bereitzustellen, sowohl für den Onlineshop als auch im Markt und per Handyapp. Durch den Einsatz von \Gls{Scrum} wird eine konstante Verbesserung der Applikation und des Arbeitsprozesses erzielt. In kontinuierlichen \Glspl{Sprint} wird die Checkout-Software auf Basis von neuen Anwendungsfällen stetig erweitert. Dieses Projekt soll dem Team als Revision dienen und Aufschlüsse über mögliche architektonische Designansätze darbieten.


\section{Motivation}

Durch den fortlaufenden Anstieg der Komplexität von Softwareprojekten \cite{Darcy.2010} haben sich gängige Designprinzipien und Architekturstile für den Entwicklungsprozess etabliert, sodass auch weiterhin die Businessanforderungen in einem zukunftssicheren Ansatz erfüllt und die Prozessabläufe jederzeit angepasst und erweitert werden können. Deshalb ist zur Gewährleistung der Langlebigkeit einer solchen Software eine flexible Grundstruktur entscheidend. Folglich ist eine sorgfältige Projektplanung und stetige Revision der Produktivanwendung relevant, um auch weiterhin einen reibungslosen Ablauf der Geschäftsprozesse zu ermöglichen. Zur Erreichung dieses Ziels verwendet die zum aktuellen Zeitpunkt bestehende Anwendung des Checkout-Teams eine Hexagonale Architektur und Domain-Driven Design.

Dieses Projekt hilft somit bei der Überprüfung der Architektur auf Verbesserungsmöglichkeiten und eventuelle Schwachstellen. Zudem existieren aufgrund des jetzigen zugrundeliegenden Aggregate-Designs gewisse Nachteile hinsichtlich der Performance und gleichzeitig stattfindenden Bearbeitung von Ressourcen. In diesem Projekt wird analysiert, ob die Performance durch eine andere Aufteilung des Datenmodells nachhaltig gesteigert werden kann. Dies dient ebenfalls als Reverenz für zukünftige Softwareprojekte, denn viele Projekte sind mit ähnlichen Problemstellungen konfrontiert.



\section{Ziele}

Aus den vorhergehenden Motivationen lassen sich folgenden Projektziele ableiten. Anhand der Analyse und Durchführung des Proof-of-Concepts wird das bestehende Softwaredesign überprüft und herausgefordert. Dadurch können konkrete Verbesserungsvorschläge an die produktive Applikation ein mögliches Fazit der Arbeit sein. Sollten sich durch ein anderes Design des Datenmodells erhebliche Vorteile ergeben, kann dies in einem Umbau der Software resultieren. Die Erkenntnisse dieser Arbeit sind ein wichtiges Ergebnis für das Team und Unternehmen, denn Projekte können auf ihrer Basis eine nachhaltigere Architektur und Datenmodell implementieren.

  
  
%% GRUNDLAGEN %% 

\chapter{Grundlagen}

Zur erfolgreichen Durchführung dieses Projekts werden Kernkompetenzen der Softwareentwicklung vorausgesetzt. Diese beschäftigen sich weitestgehend mit Softwaredesign und Architekturstilen. Um zu verstehen, wie eine Architektur die Programmierer bei der Entwicklung einer Software unterstützt, muss zunächst festgelegt werden, welche Eigenschaften der Quellcode erfüllen soll, damit dieser positive Qualitätsmerkmale widerspiegelt. Hierzu wurden gängige Designprinzipien über die Jahre festgelegt. Unter anderem die sogenannten 'SOLID'-Prinzipien, welche im nächsten Abschnitt erläutert werden. Diese tragen bei Architekturansätze miteinander zu vergleichen und bewerten.

\section{SOLID-Prinzipien}

Das weitverbreitete Akronym 'SOLID' steht hierbei für eine Ansammlung von fünf Designprinzipien, namentlich das \acrfull{SRP}, \acrfull{OCP}, \acrfull{LSP}, \acrfull{ISP} und das \acrfull{DIP}. Sie sollen sicherstellen, dass Software auch bei Expansion weiterhin testbar, anpassbar und fehlerfrei bleibt. Die grundlegenden Definitionen hinter den Begriffen lauten wie folgt:

\begin{itemize}[]
	\item \textbf{\acrlong{SRP}: } {Jede Softwarekomponente darf laut SRP maximal eine zugehörige Aufgabe erfüllen. Eine Änderung in den Anforderungen erfordert somit die Anpassung in genau einer einzelnen Komponente. Dies erhöht stark die \Gls{Kohasion} der Komponente und senkt die Wahrscheinlichkeit von unerwünschten Nebeneffekten bei Codeanpassungen.}
	\item \textbf{\acrlong{OCP}: } {Zur Sicherstellung, dass eine Änderung in einer Komponente keine Auswirkung auf eine andere besitzt, werden Komponente als 'geschlossen' gegenüber Veränderungen aber 'offen' für Erweiterungen anhand des OCPs definiert. Der erste Teil des Prinzips kann durch ein Interface realisiert werden, welches als geschlossen gilt, da die abstrakten Methoden keine Anpassungen ihrer Signatur erfahren dürfen, sonst müsste der darauf basierender Code angepasst werden. Dennoch können weiterhin Modifikationen durch das Vererben von Klassen oder die Einbindung von neuen Interfaces stattfinden. Dies wird als 'offen' im Sinne des OCPs angesehen.}
	\item \textbf{\acrlong{LSP}: } {Eine wünschenswerte Eigenschaft der Vererbung ist, dass eine Unterklasse S einer Oberklasse T die Korrektheit einer Anwendung nicht beeinflusst, wenn ein Objekt vom Typ T durch ein Objekt vom Typ S ersetzt wird. Dadurch wird die Fehleranfälligkeit bei einer Substitution im Code erheblich gesenkt und der Client kann sichergehen, dass die Funktionalität auch weiterhin den erwarteten Effekt birgt. Da das LSP sich mit der Komposition von Klassen beschäftigt, ist es für die nachfolgende Analyse vernachlässigbar.}
	\item \textbf{\acrlong{ISP}: } {Der Schnitt von Interfaces sollte so spezifisch und klein wie möglich gehalten werden, damit Clients nur Abhängigkeiten zu Funktionalitäten besitzen, welche sie wirklich benötigen. Dadurch wird die Wiederverwendbarkeit und Austauschbarkeit der Komponente gewährleistet.}
	\item \textbf{\acrlong{DIP}: } {Module sollten so unabhängig wie möglich agieren können. Dadurch wird eine erhöhte Testbarkeit und Wiederverwendbarkeit ermöglicht. Das zweiteilige DIP ist von zentraler Bedeutung für eine stabile und flexible Software. Der erste Abschnitt besagt, dass konzeptionell höherliegende Komponente nicht direkt von darunterliegenden Komponenten abhängig sein sollen, sondern die Kommunikation zwischen ihnen über eine Schnittstelle geschieht. Dies ermöglicht die Abstraktion von Funktionsweisen und löst die direkte Abhängigkeit zwischen Modulen auf. Weiterhin wird festgelegt, dass Interfaces nicht an die Implementierung gekoppelt werden sollten, sondern die Implementierung von der Abstraktion abhängen. Dadurch werden die Abhängigkeiten invertiert und ermöglicht beispielhaft die Anwendung von \Gls{DI}.}
\end{itemize}

Architekturen und Kompositionen können anhand dieser Prinzipien bewertet werden. Diese Vorgehensweise wird ebenfalls verwendet, um die nachfolgenden Architekturstile miteinander zu vergleichen. 

\section{Architekturmuster}

Eine Softwarearchitektur beschreibt die grundlegende Struktur der Module, ihre Relationen zueinander und den Kommunikationsstil unter ihnen. Die Wahl der verwendeten Architektur beeinflusst somit die komplette Applikation und ihre Qualitätsmerkmale. Das zu bevorzugende Design einer Anwendung ist stark abhängig von den Anwendungsfällen und ihren Anforderungen. 

In diesem Projekt soll ein Backend-Service erstellt werden, welcher mit den vorgelagerten Systemen über \acrshort{HTTP} und \acrshort{REST} kommuniziert. Dadurch wird die Auswahl der optimalen Architektur beschränkt, da beispielsweise Ansätze wie Model-View-Controller oder Peer-To-Peer für dieses Projektumfeld generell kaum Anwendung finden. Ein Pipe-Filter Architektur eignet sich für die Verarbeitung von einer Vielzahl an Daten, jedoch ist das Abbilden von Entscheidungsstränge und Businessrichtlinien nur umständlich verwirklichbar. Etablierte Architekturen für Backend-Software, welche die Businessprozesse als Kern der Applikation halten, müssen hingegen genauer untersucht werden. Die Schichtenarchitektur und Hexagonale Architektur erfüllen hierbei alle notwendige Bedingungen und bieten ein solides Fundament für das Projekt. Trotz ihrer ähnlichen Ziele, unterscheidet sich der Aufbau auf ersten Blick stark voneinander. In den folgenden Abschnitt werden beide Stile untersucht und anhand ihrer Tauglichkeit für eine Checkout-Software bewertet. Diese Analyse beinhaltet ebenfalls die nativ erhaltene Unterstützung der Entwickler durch die Architekturen zur Umsetzung von Designprinzipien, sodass die generelle Softwarequalität gewährleistet werden kann. 

\subsection{Schichtenarchitektur}

Durch die Einteilung der Softwarekomponenten in einzelne Schichten wird eine fundamentale Trennung der Verantwortlichkeiten und ihre Domänen erzwungen. Die Anzahl der Schichten kann je nach Anwendungsfall variieren, liegt jedoch meist zwischen drei und vier Ebenen. Die meistverbreitete Variante beinhaltet die Präsentations-, Business- und Datenzugriffsschicht. Der Kontrollfluss der Anwendung fließt hierbei stets von einer höheren Schicht in eine tiefere gelegene oder innerhalb einer Ebene zwischen einzelnen Komponenten. Ohne eine konkrete Umkehrung der Abhängigkeiten ist der Abhängigkeitsgraph gleichgerichtet zum Kontrollflussgraph. Hierbei dient Abbildung \ref{fig:Schichtenarchitektur} als eine beispielhafte Darstellung einer solchen Architektur. 

\begin{figure}[htbp]
	\centering
	\large
	\includesvg[width=0.5\textwidth]{svg/Schichtenarchitektur.svg}
	\caption{Beispielhafte Darstellung einer Drei-Schichtenarchitektur}
	\label{fig:Schichtenarchitektur}
\end{figure}

Das Ziel einer Schichtenarchitektur ist die Entkopplung der einzelnen Schichten voneinander und das Erreichen von geringen Abhängigkeiten zwischen den Komponenten. Dadurch sollen Qualitätseigenschaften wie Testbarkeit, Erweiterbarkeit und Flexibilität erhöht werden. Dank dem simplen Aufbau gewann dieser Architekturstil an großer Beliebtheit, jedoch aufgrund der fehlenden Restriktionen erhalten Entwicklern nur geringe Beihilfe zur korrekten Umsetzung des Softwaredesigns. 

Beispielsweise sind die SOLID-Prinzipien nicht oder nur minimal im Grundaufbau verankert. Das \acrlong{SRP} wird durch die Schichteneinteilung unterstützt, da Komponente zum Beispiel nicht den Zugriff auf die Datenbank und gleichzeitig Businesslogik beinhalten kann. Nichtsdestotrotz ist eine vertikale Trennung innerhalb einer Schicht nicht gegeben, daher können weiterhin Komponenten mehrere, konzeptionell verschiedene Aufgaben entgegen des SRPs erfüllen. Um die einzelnen Schichten zu entkoppelt, kann die Kommunikation zwischen den Ebenen durch Schnittstellen geschehen. Das \acrlong{OCP} soll hierbei helfen, dass Änderungen an den Schnittstellen und ihren Implementierungen die Funktionsweise, worauf tieferliegende Schichten basieren, nicht brechen. Die konzeptionelle Zuteilung dieser Interfaces ist entscheidend, um eine korrekte Anwendung des \acrlong{DIP}s zu gewährleisten. Meist wird bei sogenannten CRUD-Applikationen eine Schichtenarchitektur verwendet. \acrshort{CRUD} steht im Softwarekontext für '\textbf{Cr}eate \textbf{U}pdate \textbf{D}elete', somit sind Anwendungen gemeint, welche Daten mit geringer bis keiner Geschäftslogik erzeugen, bearbeiten und löschen. Im Kern einer solchen Software liegen die Daten selbst, dabei werden Module und die umliegende Architektur angepasst, um die Datenverarbeitung zu vereinfachen. Dadurch richten sich oft die Abhängigkeiten in einer Schichtenarchitektur von der Businessschicht zur Datenzugriffsschicht. Bei einer Anwendung, welche der Hauptbestandteil aus Businesslogik besteht, sollte hingegen die Abhängigkeiten stets zur Businessschicht fließen. Daher muss während des Entwicklungsprozesses stets die konkrete Einhaltung des DIPs beachtet werden, da entgegen der intuitiven Denkweise einer Schichtenarchitektur gearbeitet wird. Folglich bietet dieser Architekturansatz zwar einerseits einen hohen Grad an Simplizität, jedoch andererseits sind die SOLID-Prinzipien nur gering in dem Grundaufbau wiederzuerkennen.


\subsection{Hexagonale Architektur}

%Dependency Injection?

Durch weitere architektonische Einschränkungen können Entwickler zu besseren Softwaredesign gezwungen werden, ohne dabei die Implementierungsmöglichkeiten einzuengen. Dieser Denkansatz wird in der von Alistair Cockburn geprägten Hexagonalen Architektur angewandt, indem eine klare Struktur der Softwarekomposition vorgegeben wird. Hierbei existieren drei Bereiche in denen die Komponenten angesiedelt sein können, namentlich die primären Adapter, der Applikationskern und die sekundären Adapter. 

Die gesamte Kommunikation zwischen den Adaptern und dem Applikationskern findet über sogenannte Ports statt. Diese dienen als Abstraktionsschicht und sorgen für Stabilität und schützen den Kern vor Codeänderungen. Realisiert werden Ports meist durch Interfaces, welche hierarchisch dem Zentrum zugeteilt und deren Design durch diesen maßgeblich bestimmt werden. Somit erfolgt eine erzwungene Einhaltung des \emph{\acrlong{DIP}}, wodurch die Applikationslogik von externen Systemen und deren konkreten Implementierungen entkoppelt wird. Dies erhöht drastisch Qualitätsmerkmale der Anwendung, wie beispielsweise geringe Kopplung zwischen Komponenten, Wiederverwendbarkeit und Testbarkeit.

Unter den Adaptern fallen jegliche Komponenten, welche als Schnittstellen zwischen externen Systemen und der Geschäftslogik dienen. Dabei werden die primären Adapter von außerhalb angestoßen und tragen hierbei den Steuerfluss durch einen wohldefinierten Port in den Applikationskern. Zu diesen externen Systemen zählen Benutzerinterfaces, Kommandokonsolen sowie Test-Cases. Andererseits bilden alle Komponenten, welche den Steuerfluss von dem Applikationskern zu externen Systemen tragen, die Gruppe der sekundären Adapter. Hierbei entsteht der Impuls im Vergleich zu den primären Adaptern nicht außerhalb der Applikation sondern innerhalb. Die, von den sekundären Adaptern angesprochenen Systemen, können beispielsweise Datenbanken, Massage-Broker und jegliche Nachbarsysteme sein. 

Letztendlich werden alle übrigen Module im Applikationskern erschlossen. Diese beinhalten Businesslogik und sind komplett von äußeren Einflussfaktoren entkoppelt. Der beschriebene Aufbau wird in Grafik \ref{fig:HexagonaleArchitektur} veranschaulicht.

\begin{figure}[htbp]
	\centering
	\includesvg[width=0.60\textwidth]{svg/HexagonaleArchitektur.svg}
	\caption{Grundstruktur einer Hexagonalen Architektur}
	\label{fig:HexagonaleArchitektur}
\end{figure}

Das Speichern von Daten in einer hexagonalen Applikation ist ein simpler Anwendungsfall, welcher im Folgenden beispielhaft dargestellt wird. 

Ein Webclients überträgt an eine Schnittstelle des Systems Daten, wodurch er den Steuerfluss in einem sogenannten Controller initiiert. Dieser ist den primären Adaptern zugeteilt und erledigt Aufgaben, wie Authentifizierung, Datenumwandlung und erste Fehlerbehandlungen. Über einen entsprechenden Port wird der Kern mit den übergebenen Daten angesprochen. Innerhalb des Applikationszentrums werden alle business-relevanten Aufgaben erfüllt. Darunter fallen das logische Überprüfen der Daten anhand von Businessrichtlinien, Erstellen neuer Daten und die Steuerung des Entscheidungsflusses. In diesem Anwendungsfall sollen die Daten in einer Datenbank abgespeichert werden. Dementsprechend wird aus dem Anwendungskern über einen weiteren Port ein sekundäre Adapter aufgerufen, welcher für das persistieren von Daten in der Datenbank zuständig ist. \comment{Abruptes Ende. Hinleitung zur Bewertung?}

Anhand des Aufbaus einer Hexagonalen Architektur kann hinsichtlich der SOLID-Prinzipien im Vergleich zur Schichtenarchitektur folgendes Fazit formuliert werden:

\begin{itemize}[noitemsep,nolistsep,topsep=-2pt]
	\item \textbf{\acrshort{SRP}: } {Durch den Aufbau wird eine strengere konzeptionelle Trennung der Verantwortlichkeiten erzwungen. Dies wirkt sich positiv auf die Einhaltung des \acrlong{SRP}s aus.}
	\item \textbf{\acrshort{OCP} \& \acrshort{ISP}: } {Als Folge der Einführung von Ports zwischen den Applikationskern und den business-irrelevanten Komponenten ist die Anwendung der beiden Prinzipien erleichtert und teilweise vorausgesetzt. Die Applikation profitiert von erhöhter Stabilität und Kohäsion. }
	\item \textbf{\acrshort{DIP}: } {Mithilfe des \acrlong{DIP}s ist das Austauschen von Komponenten möglich, ohne dabei den Businesskern verändern zu müssen. Dies entkoppelt nicht nur den wichtigsten Bestandteil der Applikation, sondern fördert auch die Testbarkeit enorm. Durch eine native Invertierung der Abhängigkeiten bei korrekter Umsetzung der Hexagonalen Architektur gewinnt die Software vielen positiven Qualitätsmerkmalen.}
\end{itemize}


Abschließend lässt sich die Schlussfolgerung bilden, dass für eine Checkout-Software mit intensiver Businesslogik der Einsatz einer Hexagonalen Architektur zu empfehlen ist. Nicht nur ergibt sich eine natürlichere Einhaltung der SOLID-Prinzipien, sondern der Applikationskern wird ebenfalls in den Vordergrund gerückt. Anzumerken ist, dass erfahrene Entwickler jedoch ebenfalls mit einer Schichtenarchitektur ein gleiches Maß an Softwarequalität erzielen können, sofern die Designprinzipien diszipliniert eingehalten werden. 

% TODO: https://confluence.media-saturn.com/display/crosschannelcheckout/Domain+Driven+Design
\section{Domain-Driven Design}

Bei der Entwicklung von Software, welche mehr als triviale Anwendungsfälle einer CRUD-Anwendung erfüllen soll, besteht stets die Gefahr bei steigender Anzahl von Anforderungen und Änderungen zu einem sogenannten 'Big Ball of Mud' zu verschmelzen. Die bestehende Architektur wird undurchschaubar, Entstehungschancen für Bugs steigen und die Businessanforderungen finden sich überall verteilt in der Anwendung wieder. Somit kann die Wartbarkeit der Software nicht mehr gewährleistet werden und ihre Langlebigkeit ist stark eingeschränkt. Die oben analysierten Architekturstile können bei strikter Umsetzung dieses Risiko einschränken, jedoch bestimmen sie nur begrenzt wie das zugrundeliegende Datenmodell und die damit verbundenen Komponente designt werden sollen. In dem Buch \citetitle{Evans.2011} hat Eric Evans im Jahre 2003 zu diesem Zweck Domain-Driven Design entwickelt. Grundlegend wird durch diese Herangehensweise die Businessprozesse in den Vordergrund gerückt, der Problemraum in Domains eingeteilt und Richtlinien für das Design des Domainmodells festgelegt. \acrlong{DDD}, kurz \acrshort{DDD}, ist nicht gebunden an die darunterliegende Architektur oder verwendeten Technologien und an verschiedenste Einsatzgebiete anpassbar.

Bevor die Bestandteile von DDD bestimmt werden können, sollte zu Beginn eine ausführliche Umfeldanalyse der Problemebene durchgeführt werden, um festzulegen welche Verantwortungen in den zu bestimmenden Bereich fallen. Somit wird unser Problemraum als eine \emph{Domain} aufgespannt. Der Domainschnitt ist hierbei entscheidend, da basierend auf der Domain die dazugehörigen \emph{Subdomains} und ihre \emph{Bounded Contexts} bestimmt werden. Eine Subdomain bündelt die Verantwortlichkeiten und zugehörigen Anwendungsfälle in eine spezifischeren Bereich der Domain. Zur Bestimmung der Subdomains wird der Problemraum stets aus Businesssicht betrachtet und technische Aspekte werden vernachlässigt. Sollte die Domain zu groß geschnitten sein, sind dementsprechend die Subdomains ebenfalls zu umfangreich. Dadurch ist die Kohäsion der Software gefährdet und führt über den Lauf der Entwicklungsphase zu architektonischen Konflikten. Sollte eine Subdomain multiple Verantwortlichkeiten tragen, kann diese Subdomain weiter in kleinere Subdomains eingeteilt werden. Für einen Domain-Driven Ansatz ist es entscheidend die Definitionsphase gewissenhaft durchzuführen, damit eine stabile Grundlage für die Entwicklung geboten werden kann. 

Als Ausgangspunkt für die Bestimmung der Lösungsebene, werden die Subdomains in sogenannte Bounded-Contexts eingeteilt. Ein Bounded-Context kann eine oder mehrere Subdomains umfassen und bündelt ihre zugehörigen Aufgabengebiete. Wie es in der Praxis häufig der Fall ist, können Subdomains und Bounded-Context durchaus identisch sein. In jedem Bereich sollte nur ein Team agieren, um Konflikte zu vermeiden. Andernfalls kann dies ein Indiz sein, dass die Subdomains zu groß geschnitten worden sind. Jeder Bounded-Context besitzt zudem eine zugehörige Ubiquitous Language. Die Festlegung der \emph{Ubiquitous Language} stellt einen wichtigen Schritt in deinem Domain-Driven Ansatz dar. Diese definiert die Bedeutung von Begriffen, welche durch die Stakeholder und das Business verwendet werden, eindeutig. Dadurch können Missverständnisse in der Kommunikation zwischen dem Business und den Entwicklern vorgebeugt und eventuelle Inkonsistenzen aufgedeckt werden. Der größte Vorteile ergibt sich allerdings, sobald auch das Datenmodell diese Sprache widerspiegelt. Entities können Nomen darstellen, Funktionen können Verben realisieren und Aktionen können als Event verwirklicht werden. Somit sind Businessprozesse auch im Quelltext wiederzufinden. Folglich steigert dies die Verständlichkeit und Wartbarkeit der Software. Zudem lassen sich Testfälle und Anwendungsfälle leichter definieren und umsetzten. Wichtig ist, dass diese Sprache nur innerhalb eines Bounded-Context gültig ist. Beispielhaft kann der Begriff 'Kunde' in einem Onlineshop einen zivilen Endkunden, jedoch im Wareneingang eine Lieferfirma beschreiben. Daher ist bei der Kommunikation zwischen Teams unterschiedlicher Subdomains zu beachten, dass Begriffe eventuell unterschiedliche Bedeutung besitzen. 

Die  Domains, Subdomains, Bounded-Contexts und ihre Kommunikation zueinander wird durch eine Context-Map dargestellt. Diese stellt ein wichtiges Artefakt der Definitionsphase dar und kann als Tool zur Bestimmung von Verantwortlichkeiten und Einteilung neuer Anforderungen in die Domains benutzt werden. Sollte eine eindeutige Zuteilung nicht möglich sein, spricht dies für eine Entstehung eines neuen Bounded-Contexts und eventuell einer neuen Subdomain. Sowie eine Software Anpassungen erlebt, entwickelt sich die Context-Map ebenfalls stetig weiter. Zur Veranschaulichung wurde in Abbildung \ref{fig:Context-Map-Example} das Personalwesens eines Unternehmens als Domain ausgewählt und in Subdomains bzw. Bounded-Contexts aufgeteilt. Abhängig von der Unternehmensgröße und -strategie kann der Schnitt der Bounded-Contexts auch umfassender oder spezialisierter ausfallen.

%TODO: Figure als Anhang?
%TODO: Communcation between Subdomains or BoundedContexts

\begin{figure}
	\centering
	\includesvg[width=0.85\textwidth, inkscapelatex=false, height=0.85\textwidth]{svg/ExampleDomain.svg}
	\caption{Beispiel einer Context-Map anhand des Personalwesens}
	\label{fig:Context-Map-Example}
\end{figure}

Innerhalb eines Bounded-Contexts wird die grundlegende Architektur durch das zugehörige Team bestimmt. Je nach Sachverhalt des jeweiligen Kontexts kann sich diese stark von zwischen Bounded-Contexts unterscheiden. Beliebte Modellierungs- und Designstile in Verbindung mit DDD sind unter anderem Microservices, CQRS, Event-Driven Design, Schichtenarchitektur und Hexagonale Architektur. In den vorhergehenden Unterkapiteln wurden bereits die Vorzüge und Nachteile der zwei zuletzt genannten Architekturen erläutert. Auf Basis dieser Analyse wird generell für komplexere Software eine Hexagonale Architektur bevorzugt. Zudem verfolgen Domain-Driven Design und Hexagonale Architektur ähnliche Ziele, wodurch die Software natürlich an Kohäsion und Stabilität gewinnt. Im Zentrum der beiden steht das Domain-Modell, welches ohne Abhängigkeiten zu externen Modulen arbeitet. Primäre und Sekundäre Adapter sind hierzu technisch notwendige Komponente, welche durch fest definierte Ports auf den Applikationskern zugreifen können. Somit ermöglicht die Kombination aus Domain-Driven Design und Hexagonaler Architektur in Zeiten von häufigen technischen Neuheiten und komplexen Businessanforderungen weiterhin eine anpassbare, testbare und übersichtliche Software zu verwirklichen.

Auf ein solches solides Grundgerüst wird mithilfe der Kenntnisse über den Bounded-Context das Domain-Modell gesetzt. Es umfasst sowohl die Datenhaltung als auch das zugehörige Verhalten, wie zum Beispiel die Erstellung von Objekten, die Modifikation ihrer Attribute oder die dauerhafte Speicherung. Für diesen Zweck existieren in Domain-Driven Design mehrere Arten von Komponente, welche anhand ihrer Verantwortlichkeiten zugeordnet werden. Die korrekte Zuordnung der Klassen und ihrer Rollen in DDD ist entscheidend für ein skalierbares Modell, daher wird in den folgenden Unterkapiteln ein zentraler Überblick über die einzelnen Bestandteile ausgeführt.

\subsection{Value Object}

Die Value Objects bilden eine Möglichkeit zusammengehörige Daten zu gruppieren. Entscheidend ist hierbei die Frage, durch welche Eigenschaft der Zusammenschluss identifiziert wird. Die Identität eines Value Object wird alleinig durch die Gesamtheit ihrer Attribute bestimmt. Somit sind zwei Value Objects mit gleichen Werten auch identisch und miteinander ersetzt bar. 

Ein konkretes Beispiel wäre eine Klasse \emph{Preis}, welche die Attribute für Bruttobetrag, Nettobetrag und Mehrwertsteuer enthält. In den meisten Bounded-Contexts sind alleinig die konkreten Beträge von Interesse. Sollte ein Preis die gleichen Wertebelegung besitzen, gelten sie somit als identisch und austauschbar. Bei einer Aktualisierung eines Preises, kann das Objekt gelöscht und durch ein neuen Preis mit den gegenwärtigen Werten ersetzt werden. 

Aus diesen Grund gelten Value Objects als \gls{immutable}, da sie selbst keinen Werteverlauf besitzen. Dies gilt als positives Designmuster, da unveränderbare Objekte eine erhöhte Wiederverwendbarkeit genießen und unerwünschte Seiteneffekte unterdrücken. Weiterhin ist ein Resultat aus dieser Richtlinie, dass sie selbst keinen Lebenszyklus besitzen, jegliche eine Momentaufnahme des Applikationszustandes darstellen.

\subsection{Entity}

Im Gegenzug zu einem Value Object wird eine Entity nicht durch ihre Werte identifiziert, sondern behalten bei Wertanpassung die gleiche Identität bei. Folglich gelten hier zwei Entities mit einer identischer Wertbelegung als ungleich und repräsentieren unterschiedliche Objekte.

Eine Klasse \emph{Kunde} stellt einen guten Vertreter dieser Kategorie dar. In vielen Bounded-Contexts wird ein Kunde durch eine eindeutige Id ausgewiesen. Somit sind zwei Kunden mit gleichen Namen dennoch nicht die gleichen Personen. Sollte eine Person einen Namenswechsel erfahren, muss das zugehörige Objekt ebenfalls aktualisiert werden.

Daher folgt auch der Unterschied zu den Value Objects, dass eine Entity durchaus einen Lebenszyklus aufweisen. Sie werden erstellt, erhalten Anpassungen und irgendwann vom System wieder gelöscht.

In der Entity werden Businessanforderungen, welche sich auf ihre Daten beziehen, direkt implementiert. Dadurch wird eine hohe Kohäsion erzeugt und entsprechend des Information-Expert-Prinzips korrekt verankert.
%TODO: Information Expert in gls

%TODO: Hier Vergleich von Entity vs VO oder erst später?

\subsection{Aggregate}

Innerhalb des Bounded-Context ist ein Aggregate ein Verbund aus Entities und Value Objects, welches von außen als eine einzige Einheit wahrgenommen wird. Externe Komponente dürfen ein Aggregat nur durch das sogenannte Aggregate Root referenzieren und nicht direkt auf enthaltene Objekte zugreifen. Es stellt die Schnittstelle zwischen dem Aggregate und der Außenwelt dar. 

Im Bereich der Personalverwaltung ist ein mögliches Aggregat das des \emph{Mitarbeiter}s. Das Aggregat Root ist die Klasse \emph{Mitarbeiter} selbst. Diese beinhaltet Value Objects, wie \emph{Gehalt} und \emph{Abteilung}. Zu beachten ist, dass abhängig vom jeweiligen Bounded-Context zum Beispiel der Werteverlauf des Gehalt dieses Mitarbeiters vielleicht relevant ist und dementsprechend als eine Entity realisiert werden kann.

Um ein effektives Design der Aggregates zu gewährleisten, wurden einige Einschränkungen an Aggregates beschlossen. Unter anderem müssen Businessanforderungen bzw. Invarianten der enthaltenen Objekte vor und nach einer Transaktion erfüllt sein. Dadurch sind die Aggregates-Grenzen durch den minimalen Umfang der transaktionalen Konsistenz ihrer Komponente gesetzt. Zur Folge dessen, wird für jedes Select oder Update eines Datensatzes immer das komplette Aggregat aus der Datenbank geladen. Große Aggregates leiden aus diesem Grund an eingeschränkter Skalierbarkeit und Performance. Eine weitere Bedingung ist, das pro Transaktion jeweils nur ein Aggregat bearbeitet werden darf. Dies schränkt umfangreichere Aggregate durch fehlende Parallelität weiter ein. Bei Anwendung der letzteren Regel an die Mitarbeiter-Klasse, wäre es nicht möglich das Gehalt und die Abteilung durch zwei unterschiedliche Personalmitarbeiter anzupassen, da eine Transaktion der beiden Änderungen auf einen veralteten Stand operiert und zur Vermeidung eines Lost Updates zurückgerollt werden muss. 

%TODO: gls Lost Update und eventuell Select und Update

Im Falle, dass ein Anwendungsfall die Anpassung zweier Aggregates benötigt, kann dies durch eventuelle Konsistenz ermöglicht werden. Dadurch entsteht kurzzeitig ein inkonsistenter Stand der Daten, da zwei Transaktionen asynchron gestartet werden. In vielen Fällen ist ein Verzug der Konsistenz aus Sicht der Businessanforderungen akzeptabel und stellt eine gute Alternative zur Zusammenführung der beiden Aggregates dar.

%TODO: CQRS?


\subsection{Applicationservice}

Aufgaben, welche kein Domainwissen erfordert, werden in den Applicationservices realisiert. Dazu gehören das Management von Transaktionen, simple Ablaufsteuerung und Aufrufe anderer Service oder Funktionen eines Aggregate Roots. Um Nebeneffekte ausschließen zu können und Parallelität zu ermöglichen, müssen die Applicationservice ohne Zustand designt werden. Innerhalb dieser Service dürfen zudem keine Businessanforderungen enthalten sein oder Invarianten überprüft werden, dieser Aufgabenbereich wird durch die Domainservices abgedeckt.

\subsection{Domainservice}

Soweit anwendbar, werden meist alle Businessanforderungen direkt in den  zuständigen Entities oder Value Objects realisiert. Allerdings existieren Fälle, in denen keine klare Zuteilung der Aufgaben möglich ist. Dies kann beispielsweise auftreten wenn sich der Prozess über zwei oder mehr Aggregates spannt. In diesem Fall wird die Funktionalität in einem Domainservice ausgelagert. Ein weiterer Grund für die Anwendung eines Domainservices kann sein, dass die auszuführende Logik Abhängigkeiten zu anderen Services besitzt. Da das Datenmodell selbst keine Referenz auf die Domainservices haben, kann der Methodenaufruf nicht in dem Klasse selbst geschehen. Analog zu den Applicationservices werden Domainservice zustandslos implementiert.

\subsection{Factory}

Die wiederholten Erstellung von komplexen Objekten kann unnötigen Platz im Quelltext einnehmen und die Übersichtlichkeit einschränken. Dieser Effekt wird vervielfacht, sollte das Codefragment an verschiedenen Stellen auftreten. Zur Auslagerung der Objekterzeugung sind sogenannte Factories gedacht. Sie nehmen alle nötigen Daten entgegen und geben das gefragte Objekt zurück.

\subsection{Repository}

Eine Grundfunktion von allen Anwendungen stellt die Speicherung und das Laden von Daten dar. Mithilfe von Repositories wird der Datenbankzugriff ermöglicht und orchestriert. In Domain-Driven Design benötigt ein Aggregate ihr eigenes Repositorie, da sie unabhängig voneinander geladen werden müssen. Durch diese Zuordnung der Zuständigkeiten wird die konzeptionelle Abhängigkeit der Domain von den Datenbanken getrennt. Die Kommunikation mit einem Repository sollte stets über ein fest definiertes Interface geschehen, damit bei Änderungen der darunterliegenden Datenbanktechnologie der Domainkern unbetroffen bleibt. 

%TODO: Above klingt sehr stark nach Wiki: https://de.wikipedia.org/wiki/Domain-driven_Design

Mithilfe des, in diesem Kapitel erarbeiteten, Wissen wurden die Grundgedanken hinter Designprinzipien verdeutlicht und bildet somit ein solides Fundament für die Durchführung dieses Projekts. Im folgendem wird die Planungsphase des Proof-of-Concepts erläutert.



















































  
  
%% PLANUNG %% 

\chapter{Planungs- und Analysephase}

Der erfolgreiche Abschluss eines Projektes mit Domain-Driven Design erfordert die sorgfältige Analyse der Domain und des Bounded-Contexts. Basierend auf diesen Erkenntnissen kann das Domain-Modell und die Ubiquitous Language vollständig definiert werden. Besonders ist der Aggregationsschnitt von den Anwendungsfällen abhängig \cite[S. 358]{Vernon.2015}. Aus diesen Gründen wird im folgenden Kapitel der Bounded-Context umfassend untersucht. Hierbei ist die Funktionalität eines Warenkorbs und des zugehörigen Datenmodells in dem Term 'Basket' zusammengefasst, um eine einheitliche Benamung innerhalb der Ubiquitous Language einzuhalten.

\section{Ausschlaggebende Anwendungsfallbeschreibungen}


Um den Kunden im Onlineshop oder den Mitarbeitern in den Märkten eine einwandfreie Benutzererfahrung zu gewährleisten, soll die Checkout-Software alle Prozesse, vom Erstellen eines Baskets bis hin zum Kaufabschluss, verwalten können. Zu Dokumentationszwecken ist es empfehlenswert die Anwendungsfälle in Diagrammen abzubilden. Zusätzlich werden dadurch eventuelle Unklarheiten geklärt, Bedingungen an die Datenstrukturen oder den Programmfluss aufgedeckt und eine natürliche Benutzung der Ubiquitous Language etabliert. Auf Basis dieser Anwendungsfälle ist es außerdem später möglich, Artefakte des Domain-Driven Designs leichter zu definieren. Vor allem die entscheidenden Invarianten bilden sich heraus und das Datenmodell kann klarer in Aggregates unterteilt werden. Die wichtigsten Anwendungsfälle für den Proof-of-Concept sind in diesem Kapitel vereinfacht beschrieben. Bei der Reduzierung der Prozesse ist zu beachten, dass Bedingungen zwischen Datenstrukturen weiterhin unverändert bleiben, damit die möglichen Aggregationsschnitte nicht von den realisierbaren Varianten der Produktivanwendung abweichen. Dieser Abschnitt dient somit als Grundlage für das Design der Software und wird in späteren Kapiteln referenziert. Die folgenden Prozesse wurden in Zusammenarbeit mit den Teams des Onlineshops, dem Checkout-Team und zuständigen \Gls{Stakeholder}n erarbeitet. Als Darstellungsmethode wurde sich auf Aktivitätsdiagramme geeinigt, damit die Interaktionen zwischen den Systemen ebenfalls aufgezeigt werden können. Die Clients der Checkout-Software, wie der Onlineshop oder die Handy-App, sind unter dem Begriff 'Touchpoint' zusammengefasst. Falls die Graphen keine Relevanz für das leserliche Verständnis der Arbeit besitzen, sind sie dem Anhang hinzugefügt worden.

\subsection{Erstellung eines neuen, leeren Baskets}

Die Erzeugung eines Baskets kann aus verschiedenen Gründen geschehen, beispielsweise wenn ein nicht eingeloggter Kunde zum ersten Mal den Onlineshop aufruft oder ein Mitarbeiter an einer physikalischen Kasse im Markt einen Kauf abschließt. Der Erstellungsprozess besteht hauptsächlich aus dem Empfangen der Kundendaten, die Identifikation des zugehörigen Marktes mithilfe der Outlet-ID und dem permanenten Speichern des neuen Baskets. Es sind keine Invarianten zu prüfen, außer dem korrekten Format der übetragenen Nutzdaten. Dem Touchpoint wird das komplette Basket-Objekt zurückgegeben inklusive einer Basket-ID zur Referenz für spätere Zugriffe. Dieser Prozess ist in Anhang \ref{fig:SL-Basketcreation} verdeutlicht. 

\pagebreak

\subsection{Abruf eines Baskets anhand der Basket-ID}

Anhand einer ID kann ein Basket jederzeit durch den Touchpoint abgerufen werden. Die wichtigste Bedingung für diesen Anwendungsfall besagt, dass angefragte Ressourcen stets mit aktuellen Daten befüllt sein müssen. Dies stellt eine Herausforderung dar, da sich Preise und Artikeldetails mit dem Verlauf der Zeit ändern. Um das Problem möglichst performant zu lösen, werden die Informationen auf begrenzte Dauer in einem Cache zwischengespeichert. Dadurch können Aufrufe externer Systeme eingespart werden. Die genaue Speicherdauer wurde durch das verantwortliche Team festgelegt. Dies bedeutet jedoch, dass das Alter der Informationen nachvollziehbar sein muss, weshalb der Zeitstempel im Datenmodell hinterlegt wird. Auf Authentifizierung und Autorisierung wurde in dem Proof-of-Concept und den Diagrammen verzichtet, da es sich um eine rein technische Funktion handelt und keine Relevanz für die Projektumsetzung besitzt. Es verbleiben Aufgaben, wie die Suche des Baskets innerhalb der Datenbank, die De- und Serialisierung der Objekte und das Zurückgeben von Fehlermeldungen, falls beispielsweise Ressourcen nicht gefunden werden konnten. Abbildung \ref{fig:SL-Basketabruf} zeigt das zugehörige Aktivitätsdiagramm.

\begin{figure}[H]
	\centering
	\includesvg[inkscapelatex=false, width=\textwidth]{svg/AD Basketabruf.svg}
	\caption{Aktivitätsdiagramm für den Abruf eines Baskets}
	\label{fig:SL-Basketabruf}
\end{figure}

\subsection{Stornierung eines offenen Baskets}

Der Basket kann sich in verschiedenen Zuständen befinden. Es wird unterschieden zwischen 'Open', 'Frozen', 'Finalized' und 'Canceled'. Sollte beispielsweise ein Kunde beim Bezahlen an der Kasse im Markt nicht ausreichend Geld bei sich haben, muss der Basket als 'Canceled' vermerkt werden. Die tatsächliche Löschung eines Baskets ist aus rechtlichen Gründen strengstens untersagt. Für diesen Prozess ist es notwendig, vorher zu prüfen, ob der Zustand sich in 'Open' befindet, da ein eingefrorener, abgeschlossener oder bereits stornierter Basket zur Wahrung des Zustandsverlaufes nicht storniert werden kann. Dies stellt eine Invariante dar, welche in der Software sichergestellt werden muss. Im Diagramm in Anhang \ref{fig:SL-Basketstornierung} ist der Vorgang zusammenfassend abgebildet.

\pagebreak


\subsection{Hinzufügen eines Produktes anhand einer Produkt-ID}

Der Basket fungiert ebenfalls als ein Speicher einer Liste von Artikeln, sowie deren Quantität, Produktbeschreibung und ausgewählten Services bzw. Garantien. Der aufwändigste und deswegen in Abbildung \ref{fig:SL-AddProduct} dargestellte Prozess ist hierbei das Hinzufügen eines neuen Produktes zum Basket. Dabei sendet der Client lediglich die zugehörige Produkt-ID, weswegen externe Systeme von der Checkout-Software aufgerufen werden müssen. Dazu gehört die Product-API, welche alle notwendigen Produktdetails liefert. Der Artikelpreis selbst ist hierbei nicht in den Produktinformationen zu finden, da dieser von Markt zu Markt unterschiedlich sein kann. Daher wird eine weitere API-Abfrage benötigt, welche zu einer Produkt-ID und Outlet-ID den zugehörigen, aktuellen Preis zurückschickt. Diese zwei Schnittstellen müssen ebenfalls in anderen Anwendungsfällen aufgerufen werden, sofern die zwischengespeicherten Preis- bzw. Artikelinformationen im Cache nicht mehr gültig sind. Zusätzlich folgt eine Validierung des Datenmodells auf verschiedene Parameter. Unter anderem darf die Gesamtanzahl der enthaltenen Artikel keinen festen Wert überschreiten. Der aktualisierte Basket wird beim erfolgreichen Abschluss der Operation dem Touchpoint zurückgegeben.

\begin{figure}[H]
	\centering
	\includesvg[inkscapelatex=false, width=\textwidth]{svg/AD AddProduct.svg}
	\caption{Aktivitätsdiagramm für das Hinzufügen eines Produktes }
	\label{fig:SL-AddProduct}
\end{figure}

\subsection{Hinzufügen einer Zahlungsmethode}

Eine weitere essentielle Funktion ist das Hinzufügen von Zahlungsmethoden, damit der Kauf erfolgreich initiiert werden kann. Da es sich nur um das ungeprüfte Anhängen der Bezahlinformationen handelt, müssen keine strengen Validierungen vorgenommen werden, weil diese Aufgabe durch ein externes System in einem späteren Schritt des Checkouts erledigt wird. Dennoch sind logische Überprüfungen durchzuführen, wie zum Beispiel, dass der Basket nicht leer oder bereits bezahlt ist. Ebenfalls findet eine Neuberechnung aller Bezahlinformationen statt. Analog zu den vorgehenden Fällen wurde ein Aktivitätsdiagramm in Anhang \ref{fig:SL-PutBezahlmethode} designt.

\pagebreak

\subsection{Aktualisieren der Checkout Daten des Baskets}

Ein Basket besitzt eine große Menge an Attributen. Einige dieser werden implizit durch einen Prozess innerhalb der Software gesetzt, andere durch Empfangen der Daten von einem externen System. Beim sogenannten 'Checkout-Prozess' werden Informationen, wie Kundendaten, Zahlungsmethode und Zustellungsart, von den Touchpoints an die API gesendet. Da die Daten innerhalb eines einzelnen Verfahrens durch das vorgelagerte System gesammelt werden, bietet es sich an, die korrespondierende Schnittstelle so zu designen, dass alle betroffenen Attribute gleichzeitig angepasst werden können. Eine Überprüfung der Daten erfolgt in diesem Schritt nicht, mit der Ausnahme, dass die gewählte Zustellungsart, in dem Bounded-Context 'Fulfillment' genannt, für den ausgewählten Basket verfügbar sein muss. Zudem ist es notwendig eine Neuberechnung der Zahlungsmethoden durchzuführen, falls eine neue Zahlart hinterlegt worden ist. Diese Bedingungen und der dazugehörige Prozess sind in Anhang \ref{fig:SL-Checkoutdata} veranschaulicht worden.


\subsection{Initiierung des Bezahlprozesses und Einfrieren des Baskets}

Nachdem die Bearbeitung des Baskets abgeschlossen ist, kann der Bezahlprozess gestartet werden. Hierbei muss das Datenmodell einen konsistenten und validen Stand besitzen. Sollte dies der Fall sein, wird der Basket in den Status 'Freeze' gestellt und jegliche weitere Datenänderungen verhindert. Der Bezahlprozess wird von einer externen Software abgewickelt und durch eine Referenz im Basket verlinkt. Eine vereinfachte Darstellung des Prozesses bietet Abbildung \ref{fig:SL-InitPayment}.

\begin{figure}[H]
	\centering
	\includesvg[inkscapelatex=false, width=\textwidth]{svg/AD InitPayment.svg}
	\caption{Aktivitätsdiagramm für das Initiieren des Bezahlvorgangs }
	\label{fig:SL-InitPayment}
\end{figure}

\pagebreak

\subsection{Ausführung des Bezahlprozesses und Finalisierung des Baskets}

Als letzter, höchst relevanter Anwendungsfall gilt der Abschluss des Bezahlvorgangs, visualisiert in Abbildung \ref{fig:SL-ExecPayment}. Die Checkout-Software dient hierbei als Proxy zwischen den Touchpoints und der Payment-API. Sofern die Bezahlung erfolgreich war, wird der Status des Baskets auf 'Finalized' gestellt. Zugleich wird eine Bestellung durch die Order-API angelegt und im Basket durch eine Referenz verlinkt.

\begin{figure}[H]
	\centering
	\includesvg[inkscapelatex=false, width=\textwidth]{svg/AD ExecPayment.svg}
	\caption{Aktivitätsdiagramm für das Ausführen des Bezahlvorgangs }
	\label{fig:SL-ExecPayment}
\end{figure}

\subsection{Aus Anwendungsfällen resultierende API-Schnittstellen }

Anhand der Anwendungsfälle ist es möglich eine klare Schnittstellendefinition für die Checkout-Applikation zu erstellen. Hierbei beinhaltet diese alle benötigten Operationen zum erfolgreichen Bewältigen der Anforderungen aus Sicht der Touchpoints. Die Kommunikation der Systeme geschieht über eine \acrshort{REST}-API und somit auf Basis des \acrshort{HTTP}-Protokolls. In Anhang \ref{fig:REST-API} sind alle Endpunkte enthalten, mitsamt ihrer HTTP-Methode, Parameter und gesendeten bzw. empfangenen Datenstrukturen.

\pagebreak

\section{Projektumfeld und technologische Vorschläge}

Die komplette Systemumgebung von MediaMarktSaturn ist eine komplexe Architektur mit zahlreichen Abhängigkeiten zwischen Teams und ihren betreuten Applikationen. Es ist unmöglich ein solches Konstrukt aufzubauen ohne die Kommunikation der einzelnen Systeme untereinander zu definieren. Als Leitfaden für das Projektumfeld dienen die Anwendungsfälle des vorgehenden Unterkapitels. In dem vereinfachten Checkout-Prozess werden sechs verschiedene Schnittstellen aufgerufen. Damit plötzliche Systemänderungen keine Auswirkung auf die Funktionsweise der abhängigen Clients haben, wird eine verpflichtende API-Vereinbarung beschlossen, welche die benötigten Informationen, mögliche Fehlerfälle und die zurückgelieferten Daten festlegt. Der Proof-of-Concept orientiert sich an diesen Vereinbarungen, erleichtert allerdings die Kommunikationsbedingungen, um unnötigen \emph{\Gls{Boilerplate-Code}} zu unterdrücken. Die Abbildung \ref{fig:ContextDiagramm} stellt ein Context-Diagramm der Umgebung dar. 

\begin{figure}[htbp]
	\vspace{0.5cm}
	\centering
	\includesvg[inkscapelatex=false, width=0.8\textwidth]{svg/ContextDiagramm.svg}
	\caption{Context Diagramm der produktiven Checkout-Umgebung}
	\label{fig:ContextDiagramm}
	\vspace{0.5cm}
\end{figure}

Zusätzlich bestehen noch firmen- bzw. teaminterne Vorbedingungen. Wo sinnvoll anwendbar, wird in der Firma Java als Entwicklungssprache angewandt. Über die letzten Jahre gewann Kotlin an Beliebtheit und wird seitdem ebenfalls bei MediaMarktSaturn eingesetzt. Die aktuelle Live-Umgebung nutzt zu dem jetzigen Stand noch Java, jedoch werden Codeanpassungen zukünftig in Kotlin vorgenommen, um eine langsame Migration zu gewährleisten. Aus diesem Grund wird ebenfalls der Proof-of-Concept in Kotlin umgesetzt. Zudem ist die Technologie der systemübergreifenden Kommunikation auf REST-APIs festgelegt. Dies kommt mit einigen Einschränkungen und muss in der Entwicklung der primären und sekundären Adapter beachtet werden. Die Auswahl der Datenbank ist grundsätzlich nicht vorgegeben. In einem späteren Kapitel wird die Verwendung einer MongoDB für das Projekt begründet.

  
  
%% PROOF OF CONCEPT %% 

\chapter{Erstellung eines Proof of Concepts}

Durch die Schaffung eines grundlegenden Verständnis für Designprinzipien, Hexagonaler Architektur und Domain-Driven Design kann auf Basis der vorgehenden Analysen ein Proof of Concept entwickelt werden. Hierzu wird weiterhin das typisches Vorgehen von DDD verfolgt und zunächst die Domain und Ubiquitous Language definiert, gefolgt vom Erstellen des zentralen Domain-Modell. 

Die Wahl der Programmiersprache ist aufgrund der Teamexpertisen und dem vorliegenden Anwendungsfall auf Java oder Kotlin beschränkt. Kotlin ist eine, auf Java basierende, Programmiersprache, welche in ihrem Design einen hohen Fokus auf Lesbarkeit und Flexibilität legt. Aufgrund einer Vielzahl von Vorteile, welche durch den Einsatz von Kotlin erzielt werden können, wurde diese als die zu verwendende Sprache festgelegt. Zugleich wird nach Projektabschluss jegliche Codeanpassung an der produktiven Checkout-Software in Kotlin vorgenommen und somit ein langsame Migration des Sourcecodes erreicht.

%TODO: Wie Einbindung des Source Codes
%TODO: Bild der Architektur
%TODO: Domain oder Domäne
%TODO: Programmiersprache, REST, MongoDB in Umfeldanalyse oder hier ein extra Unterkapitel?

\section{Festlegung und Analyse der Domain}
\blindtext

\section{Abgrenzung des Bounded Contexts}
\blindtext

\section{Festlegen einer Ubiquitous Language}

In der Kommunikation zwischen den Business und Entwicklern kann es oft zu Missverständnisse kommen. Womöglich werden Informationen, Einschränkungen oder Prozesse ausgelassen oder für selbstverständlich erachtet. Durch die klare Definition von gemeinsam verwendeten Begriffen und ihren Bedeutungen wird implizit notwendiges Wissen über die Domain und ihre Prozesse geschaffen. Viele dieser Fachbegriffe können für notwendige Anwendungsfälle wiederverwendet werden und machen die Personen, welche die Businessanforderungen umsetzten sollen, mit der Domain vertraut. Da mit geringem Domainwissen die Korrektheit der Software gefährdet wird, stellt die Ubiquitous Language in Domain-Driven Design ein wichtiger Teilaspekt dar.

In Zusammenarbeit mit den Lead-Developer und \Gls{Product Owner} des Teams wird im folgenden Abschnitt die Ubiquitous Language definiert, um ein solches Verständnis über den Checkout Bounded-Context zu gewährleisten. Hierbei wurde sich auf die, für dieses Projekt, relevanten Terme beschränkt. Eingeklammerte Wörter beschreiben Synonyme zu dem vorangestellten Ausdruck.

{\large \textbf{Domain-Modell:}}
\begin{itemize}[noitemsep,nolistsep]
	\item \textbf{Basket (Warenkorb): } {Die zentrale Datenklasse, welche den Warenkorb darstellt.}
	\item \textbf{Customer (Kunde): } {Ein Endkunde des Onlineshops oder im Markt. Kann eine zivile Person sein oder eine Firma.} %TODO: Kann er beides sein, wie funktioniert das aus business sicht?
	\item \textbf{Product (Artikel, Ware): } {Ein Artikel aus dem Warenbestand, welcher zu Verkauf steht. Kann ebenfalls für eine Gruppierung von mehreren Artikeln stehen.}
	\item \textbf{Outlet: } {Repräsentiert einen Markt oder den länderbezogenen Onlineshop, welche durch eine einzigartige Outlet Nummer referenziert werden können.}
	\item \textbf{BasketItemId: } {Eine, innerhalb eines Baskets, eindeutige Referenz auf ein Artikel des Warenkorbs. Wird aus technischen Gründen benötigt, um beispielsweise die Quantität eines Artikels anzupassen oder ihn zu entfernen.}
	\item \textbf{Net Amount (Nettobetrag): } {Ein Nettobetrag mit dazugehöriger Währung.} %TODO
	\item \textbf{VAT (Steuersatz): } {Der Steuersatz zu einem bestimmten Nettobetrag.}
	\item \textbf{Gross Amount (Bruttobetrag): } {Der Bruttobetrag eines Preises errechnet aus dem Steuersatz und einem Nettobetrag. Die Währung gleicht die des Nettobetrags.}
	\item \textbf{Fulfillment: } {Zustellungsart der Waren seitens der Firma.}
	\begin{itemize}[noitemsep,nolistsep]
		\item \textit{Pickup: } {Warenabholung in einem ausgewählten Markt durch den Kunden. Nur möglich sofern Artikel im Markt auf Lager ist.}
		\item \textit{Delivery: } {Zustellung der Ware an den Kunden durch einen Vertragspartner.}
		\item \textit{Packstation: } {Lieferung der Ware an eine ausgewählte Packstation durch einen Vertragspartner.}
	\end{itemize}
\end{itemize}
\vspace{0.4em}
%TODO: Maybe add PriceAdjustment

{\large \textbf{Prozesse:}}
\begin{itemize}[noitemsep,nolistsep]
	\item \textbf{Basket Validation: } {Durchführung einer Validierung des Baskets auf Inkonsistenzen oder fehlenden, notwendigen Werten.}
	\item \textbf{Basket Finalization: } {Überprüfung der Korrektheit des Baskets durch die Basket Validation, Sperrung des Baskets vor weiteren Abänderungen, Einleitung des Zahlungsprozesses und Reservierung der Produkte im Basket. }
	\item \textbf{Basket Cancellation: } {Stornieren des zugehörigen Baskets. Kann nur auf einen nicht geschlossenen oder stornierten Basket ausgeführt werden. Nach Cancellation dürfen keine weiteren Änderungen an den Basket durchgeführt werden. }
	\item \textbf{Basket Creation: } {Explizite oder Implizite Erstellung eines neuen Baskets. Findet automatisch statt sofern noch kein Basket für den Customer existiert. Dies ist ebenfalls der Fall nach einer Basket Finalization.}
	\item \textbf{Basket Calculation: } {Das Berechnen der Geldbeträge des Baskets. Beinhaltet die Auswertung aller Products mitsamt ihren Net Amounts und VATs. Der resultierende Betrag aus unterschiedlichen VATs muss weiterhin aus rechtlichen Gründen einzeln verwiesen werden können.}
	\item \textbf{High Volume Ordering: } {Die Bestellung von Artikeln in hoher Stückzahl. Aufgrund von Businessanforderungen soll es nur begrenzt möglich sein, dass ein Kunde in einem Basket oder wiederholt das gleiche Produkt mehrfach kauft.}
\end{itemize}
\vspace{0.4em}

%TODO: Basket Finalization oder finalize Basket?
%TODO: Lieber Wortwiederholungen oder Abwechslung
%TODO: Stil der Sätze... Punkt nach unvollständigen Satz?
%TODO: Technische Begriffe in Language? "BasketItemId" etc.
%TODO: Doppelt gemoppelt der obere Absatz mit den unteren?

Im Verlaufe der Definitionsphase der Ubiquitous Language wurden die Prozesse näher beleuchtet, Benamungen von Datenobjekten aufdeckt und Businessanforderungen vorgegeben. Ein gutes Datenmodell spiegelt hierbei ebenfalls die Sprache des Bounded-Contexts wieder, daher wird auf Basis dieses Unterkapitels und den vorgehenden Analysen anschließend das Datenmodell designt. 

\section{Definition der Value Objects}
Aufgrund der Simplizität und klaren Zuteilung vom Modell zu Value Objects lassen sich diese als leichtes bestimmen. Value Objects kapseln sowohl Daten als auch das dazugehörige Verhalten. Wobei ihre Funktionalitäten keine Änderungen an den beinhalteten Attributen durchführen dürfen, da sie als immutable gelten. Daher wird, anstatt einer direkten Wertanpassung, in der Regel eine Kopie des aktuellen Objektes mitsamt der aktualisierten Werte erstellt und zurückgegeben. Zudem gilt diese Kopie als ein neues Value Object, dank der Eigenschaft, dass die Identität eines Value Objects alleinig von ihren Attributen und den zugeteilten Werten abhängt. Anfangs sollte jede Datenstruktur des Domain-Modells als ein Value Object definiert und erst nach gründlicher Überlegung, falls die Notwendigkeit besteht, zu einer Entity umgeschrieben werden.

Anfangs wird das Daten-Modell mithilfe dem erlangten Domainwissen grundlegend aufgebaut. Die Ubiquitous Language unterstützt bei der richtigen Benamung der Klassen. Der Basket stellt den Ausgangspunkt der Modellierung dar. Folgende Attribute werden benötigt:

\domainModell{Basket}{
	\item \textbf{OutletId: } {Zur Referenz in welchem Markt oder Onlineshop der Warenkorb angelegt wurde. Unerlässlich für die Bestimmung von unter anderem Lagerbeständen, Lieferzeiten, Fulfillment-Optionen und Versandkosten. Wir benötigen nur die Id des Outlets und nicht alle Daten, da sie für unsere Kommunikationspartner nicht relevant sind.}
	\item \textbf{Customer: } {Speichert Kundendaten.} %TODO: Warum benötigt? Wird weitergegeben an COS?
	\item \textbf{BasketItems: } {Liste an allen Artikeln aktuell im Warenkorb.} %TODO: Warum nicht Product hier erklären
	\item \textbf{CalculationResult: } {Beinhaltet die berechneten Werte des Warenkorbs, wie Nettobetrag, Bruttobetrag und VAT. Die Speicherung dieser Werte wäre technisch nicht notwendig, spart aber an Rechenzeit ein, da nicht bei jeder Abfrage des Basket dieser Wert neu berechnet werden muss.}
	\item \textbf{Status: } {Repräsentiert, ob der Warenkorb als offen, canceled oder finalized gilt. Wichtig für Businessprozesse und Validierung von Änderungen am Basket.} %TODO: Umbenennung im Source code von COMPLETED zu FINALIZED
}

Anschließend werden die, im Basket enthaltenen, Klassen ebenfalls mit der gleichen Vorgehensweise definiert, sofern sie mehr als nur einfache Strings oder \Glspl{Enumeration} darstellen:

\domainModell{Customer}{
	\item \textbf{Name: } {Enthält den Vor- und Nachnamen als eigenes Datenkonstrukt.}
	\subDomainModell{ %TODO: Untergliedern?
		\item \textit{FirstName: } {Vorname des Kunden.} %TODO: Punktsetzung?
		\item \textit{LastName: } {Nachname des Kunden.}
	}
	\item \textbf{Email: } {Email des Kunden.} %TODO: Warum benötigt?
	\item \textbf{CustomerTaxId: } {Die Steuernummer des Kunden. Relevant aus Sicht der Rechnungsabwicklung und für den Ausdruck der Rechnung, da diese ausgewiesen werden muss} %TODO: ///    Muss TaxId ausgewiesen werden?
	\item \textbf{Type: } {Bestimmt ob Kunde als B2C oder B2B gilt.} %TODO: gls? /// Wie wird bestimmt?
	\item \textbf{CompanyName: } {Firmenname des Kunden. Kann optional angegeben werden oder ist verpflichtend bei einem B2B-Kunden} %TODO: Korrekt?
	\item \textbf{CompanyTaxId: } {Steuernummer zugehörig zur Firma des Kunden, falls eine Firma angegeben wurde.} %TODO: Korrekt? Wofür? Warum CustomerTaxId und CompanyTaxId? TaxId in Language?
}

\domainModell{BasketItem}{
	\item \textbf{Id: } {Eindeutige Nummer des Items in diesem Basket.}
	\item \textbf{Product: } {Beinhaltet alle Produktinformationen, welche durch die Software benötigt werden.}
	\item \textbf{Price: } {Aktueller Preis des zugehörigen Products. Kann sich zeitlich ändern, muss daher durch eine Businessfunktion aktualisiert werden, jedoch beeinflusst dies nicht die Modellierung des Datenobjekts. } %TODO: Letzten Satz so lassen oder weglassen? Anderes Wort für Relevanz?
}

\domainModell{Product}{
	\item \textbf{Id: } {Eindeutige Referenz für dieses Product im externen System.}
	\item \textbf{Name: } {Textuelle Produktbezeichnung des Products.}
	\item \textbf{Vat: } {Mehrwertsteuerinformationen des Products.}
}

\domainModell{Vat}{
	\item \textbf{Sign: } {}
	\item \textbf{Rate: } {Prozentualer Wert der Mehrwertsteuer.} %TODO: Bedeutung?
}

\domainModell{Price}{
	\item \textbf{GrossAmount: } {Bruttobetrag mit Währung.} %TODO: In POC nur Gross als Price?
	\item \textbf{NetAmount: } {Nettobetrag mit Währung.} %TODO: Beschreibung weglassen wenn gleich mit Language?
	\item \textbf{VatAmount: } {Mehrwertsteuerklasse des Preises. Benötigt für die Berechnung des Preises.}
}

\domainModell{VatAmount}{
	\item \textbf{Sign: } {}
	\item \textbf{Rate: } {Prozentualer Wert der Mehrwertsteuer.}
	item \textbf{Amount: } {Berechneter Betrag der Mehrwertsteuer zugehörig zu einem Bruttobetrag.}
}

\domainModell{CalculationResult}{
	\item \textbf{GrossAmount: } {Bruttobetrag mit Währung.}
	\item \textbf{NetAmount: } {Nettobetrag mit Währung.}
	\item \textbf{VatAmounts: } {Eine zusammengebautes Set aus VatAmounts der Preise der BasketItems. Benötigt, da Vats mit unterschiedlichen Prozentbeträgen rechtlich nicht kombiniert werden dürfen.}
}

%TODO: Neuanordnung der Attribute von Customer im Source Code

%TODO: Datentyp hinter name?
%TODO: Definition doppelt?
%TODO: Zu oft Value Objects
%TODO: Klassendiagram
%TODO: Wie darstellen?

\section{Bestimmung der Entities anhand ihrer Identität}

Anhand der vorgehenden Sektion ist das Datenmodell nun vollständig. Jedoch besteht weiterhin die Frage, ob die jeweiligen Klassen eine eigene Identität besitzen und somit als Entity designt werden müssen. Es existiert in DDD kein objektive Verfahren zur Bestimmung der Entities, da je nach Bounded-Context Datengruppierungen unterschiedliche Eigenschaften besitzen. Als Hilfestellung für diese Entscheidung existieren folgende grundlegende Unterscheidungsmerkmale und Richtlinien:

\begin{tabular}{ | p{0.16\textwidth} | p{0.27\textwidth} | p{0.47\textwidth} | } 
	\hline
	& {\textbf{\centering Value Object} } & \textbf{\centering Entity} \\ 
	\hline
	Identität & 
	Summe der Werte der Objekte. Objekte mit gleichen Werten besitzten gleiche Identität & 
	Bestimmt anhand eines Identifikator, zum Beispiel einer Datenbank-Id. Objekte mit gleichen Werten sind ungleich, außer ihre Identifikatoren sind identisch. \\ 
	\hline
	Lebenszyklus & 
	Stellt nur eine Momentaufnahme des Applikationszustands dar. &
	Werden zu einem bestimmten Zeitpunkt erstellt, bearbeitet, gespeichert oder gelöscht und besitzen somit einen Verlauf ihrer Wertänderungen.  \\ 
	\hline
	Veränderbarkeit & 
	Durch einen fehlenden Lebenszyklus gelten Value Objects als immutable. &
	Aufgrund ihrer Eigenschaften sind Entities mutable.  \\ 
	\hline
	Abhängigkeit & 
	Können nur als Unterobjekt einer oder mehrerer Entities existieren. &
	Um einen eigenen Lebenszyklus zu besitzen, können sie unabhängig von anderen Objekten erstellt werden.  \\ 
	\hline
	Zugriffsmethode & 
	Auf Daten und Funktionen wird mithilfe einer Entität zugegriffen. &
	Können als Aggregate Root direkt Zugriff erfahren. \\ 
	\hline
\end{tabular}

Anhand dieser Eigenschaften können die Value Objects untersucht und daraufhin alle Entities bestimmt werden:

%TODO: Subsection oder itemize, oder einfach fette überschiften und dann text?
\domainModellWithoutTitle{
	\item \textbf{Basket: } {Als zentrales Datenobjekt besitzt ein Basket zur eindeutigen Identifikation eine Id, dies spricht stark für eine Entity. Zusätzlich bestimmen nicht die enthaltenen Attribute wie Products oder der zugehörige Kunde die Identität des Baskets. Aufgrund der geforderten Anwendungsfälle entsteht zugleich ein Lebenszyklus für die Instanzen eines Baskets.}
	\item \textbf{Customer: } {In Bounded-Contexts, welche mit den Kundendaten operieren, kann diese Klasse durchaus eine Entity darstellen. In unserer Domain finden keine Operationen auf diesen Informationen statt und die vorangestellten Systeme senden bei Änderungen der Kundendaten diese zu. Folglich besitzen sie keinen eigenen Lebenszyklus, werden ebenfalls nicht separat gespeichert und kann weiterhin als Value Object designt werden.}
	\item \textbf{BasketItem: } {}
	\item \textbf{Product: } {}
	\item \textbf{Price: } {}
	\item \textbf{CalculationResult: } {Als Datenstruktur, welche bei jeder Neuberechnung aktualisiert wird, könnte die Eigenschaft eines Lebenszyklus erfüllt sein. Jedoch stellt die Klasse einzig ein Zwischenspeicher der Ergebnisse dar und ohne den Kontext eines darüberlegenden, zugehörigen Objekt besitzen diese Daten keine Aussagekraft. Weiterhin sind die gleichen Ergebnisse im Sinne der Identität äquivalent. }
}

%TODO: Wieder doppelt? Schon in DDD Kapitel? Als Tabelle oder als itemize?
%TODO: Ref to table






  
  \chapter{Design der Aggregates anhand der Anwendungsfälle}
%TODO: Zeigen des Bestellprozesses anhand von MediamarktSaturn

Ein Domainmodell kann unterschiedlich gestaltet werden. Ebenfalls betrifft dies den Schnitt der Aggregates. Verschiedene Schnitte genießen verschiedene Vor- und Nachteile. In diesem Kapitel werden mehrere mögliche Aufteilungen untersucht und durch verschiedene Kriterien bewertet. Diese Kriterien umfassen unter anderem Performance, Komplexität, Konsistenz, Parallelität, Relevanz und Anwendbarkeit im Bezug auf den Bounded Contexts des Checkouts.

\section{Ein zusammengehöriges Basket-Aggregate als initiales Design}

Das Design eines großen Aggregats fällt Entwicklern meist natürlich leichter. Eine einzelne zusammenhängende Struktur, durch welche Daten einfach bearbeitet und Invarianten überprüft werden können. Daher hält sich die Komplexität bei diesem Aggregationsschnitt weitestgehend in Grenzen. Vor allem müssen kaum Überlegungen über transaktionale Konsistenz getroffen werden, da die Grenzen der Transaktion all umspannend ist. Die erste Variante des Proof-Of-Concepts wurde mit diesem Design im Gedanken entwickelt und stellt das Grundgerüst für alle folgenden Konzepte dar. Zur vereinfachten Referenz auf die einzelnen Umsetzungsvarianten wird das Design, welches einen großen Basket zugrundeliegend hat, als '\emph{Variante A}' betitelt.
%TODO: Variante A als Bezeichnung ok???


\subsection{Performance von unterschiedlich großen Aggregates im Vergleich}

%TODO: Lazy loading
Als Konsequenz eines großen Aggregates ergibt sich, dass immer das ganze Aggregate geladen werden muss, da auf darunterliegende Datenstrukturen nur durch das Aggregate Root zugegriffen werden darf. Für die Aktualisierung einer Entity tief im Aggregat müssen alle anderen Daten ebenfalls geladen werden. Die Auswirkungen dieser Aussage kann auf Datenbankebene anhand von Lazy-Loading oder durch Einsatz einer dokumentbasierten Datenbank eingeschränkt werden. Generell gilt, dass große Aggregates aus Sicht der Performance langsamer arbeiten als kleinere. Diese Aussage ist allerdings mit Vorsicht zu genießen und kann je nach Anwendungsgebiet sogar gegensätzlich ausfallen. In diesem Unterkapitel wird die diese Richtlinie anhand unseres Bounded Contexts analysiert.

Wie oben erwähnt, ermöglicht ein kleinerer Aggregationsschnitt das unabhängige laden und bearbeiten der Aggregates. In Betrachtung der Anwendungsfälle können die Operationen in drei große Gruppen umfasst werden. Anfragen, welche Value Objects im Basket bearbeiten stellt die erste Ansammlung dar. Diese würden unnötige Daten laden, da nur der Basket mit möglichst wenig Value Objects relevant ist. Ebenfalls gilt dies für die Anfragen, welche BasketItems oder den PaymentProcess anpassen, weil sie den ganzen Basket laden anstatt nur die wirklich benötigten Daten.

Wird die \emph{Variante A} unterteilt in Basket-, Payment und BasketItem-Aggregate (\emph{Variante B}) können diese somit unabhängig voneinander geladen und bearbeitet werden. Ohne Beachtung, welche Auswirkungen dieses neu überlegte Design auf andere Faktoren hat, ist eine Performanceverbesserung anhand ersten Überlegungen zu erwarten. Die vermeintlich gewonnene Performance wird allerdings aufgrund mehrerer Umständen verringert. Eine kurze Übersicht über diesen neuen Aggregationsschnitt bietet Abbildung \ref{fig:VarB}.

\begin{figure}[htbp]
	\centering
	\includesvg[inkscapelatex=false, width=0.80\textwidth]{svg/VarB.svg}
	\caption{Aggregationsschnitt der Variante B}
	\label{fig:VarB}
\end{figure}

\subsubsection{Invarianten über Aggregatsgrenzen hinweg}

Wird dem Warenkorb ein neuer Artikel hinzugefügt hat dies zu Folge, dass nicht nur ein neues Item angelegt, sondern auch der Baskets neu berechnet werden muss. Sofern eine Aggregatstrennung zwischen Basket und BasketItem vorliegt, muss dennoch der ganze Basket geladen werden, andernfalls ist es nicht möglich den Gesamtpreis zu ermittelt. Innerhalb eines Checkout-Contexts existieren viele dieser Businessanforderungen, welche mehrere Objekte benötigen, um angewandt zu werden. Zudem ist es gegen die grundlegenden Bestimmungen eines Aggregates, dass eine Anfrage bzw. Transaktion zwei oder mehrere Aggregate bearbeiten. Dieser Anwendungsfall hätte eine Speicherung eines neuen BasketItems und die Aktualisierung des CalculationResults des Baskets zufolge. Schlussfolgernd wäre eine solche Operation nicht erlaubt ohne Verwendung von eventueller Konsistenz. Um genauer zu bestimmen, welche tatsächliche Performance erzielt werden kann durch kleinere Aggregates am konkreten Design der \emph{Variante B} müssen die Anzahl der tatsächlich unabhängig abschließbaren Operationen untersucht werden. Die vorher festgelegten Swimlane Diagramme helfen bei dieser Aufgabe.

Ein Abruf eines Baskets kommt einher mit dem Laden aller anderen Aggregates, da die Touchpoints den gesamten Basket benötigen. Dies gilt ebenfalls für alle anderen Anwendungsfälle, welche als Antwort auf ihre Anfrage einen kompletten Basket erwarten. Das Stornieren des Baskets und Setzten der Checkout-Daten ist positiv von dieser Anpassung betroffen. Jedoch ist die Initiierung und Durchführung des Bezahlvorgangs nicht komplett unabhängig abschließbar, da im Nachhinein der Status des Baskets angepasst werden muss. Dadurch entsteht eine Abhängigkeit, welche im Nachhinein erfüllt werden muss und somit zur eventuellen Konsistenz führt. Die größte Anzahl an API-Anfragen beziehen sich auf das Hinzufügen und Bearbeiten von BasketItems und die Verwaltung des Bezahlvorgangs. Letztendlich sind nur eine Bruchzahl der Anwendungsfälle durch einen Umbau der Architektur von Variante A auf Variante B wirklich positiv betroffen.

\subsubsection{Einfluss des Datenbankmodells auf den Aggregationsschnitt}

Das Design einer Software soll stets so weit wie möglich unabhängig von verwendeten Technologien sein. Technologien entwickeln sich weiter, werden durch neuer ersetzt und bringen unnötige Abhängigkeiten in den Quelltext. Theoretisch hat somit eine Beeinflussung der Architektur durch eine externe Komponente einen negativen Effekt auf die Qualität der Software und ihre Wartbarkeit. Dennoch kann in der Praxis dieser Gedanke durchaus Vorteile bergen, welche dieses Vorgehen gerechtfertigt. Folglich wird der Aggregatsschnitt aus Sicht der Datenbank bewertet.

Laut Definition erhält jedes Aggregate seine eigene Tabelle in der darunterliegenden Datenbank. Daher existieren in \emph{Variante B} mindestens die Tabelle für den Basket, BasketItems und PaymentProcess. In \emph{Variante A} kann der PaymentProcess in die gleiche Tabelle wie der Basket geschrieben werden, da eine Eins-zu-Eins Relation vorliegt. Obwohl es nur ein Aggregate gibt, sind dennoch zwei Tabellen für Basket und BasketItem notwendig, da in einer Relationalen Datenbank eine Eins-zu-N Beziehung nicht anders dargestellt werden kann. Demnach muss beim Abruf eines kompletten Baskets in \emph{Variante A} N Joins durchgeführt werden, wobei \emph{Variante B} N+1 Joins benötigt, wobei N die Anzahl der BasketItems entspricht. Dies resultiert in eine schlechtere Performance der \emph{Variante B} zumindest in Bezug der Datenbankabfrage. Wie bereits erwähnt, erwarten eine Vielzahl von Anfragen alle Daten des Baskets zurück und dadurch erhalten wir eine erhöhte Bearbeitungszeit der Anfragen.

%TODO: Performaceanalyse im POC
%TODO: Collection
Sollte die verwendete Datenbank jedoch einen dokumentbasierten Ansatz verfolgen, gilt vorheriger Absatz nur noch bedingt. Hierbei benötigt \emph{Variante B} weiterhin drei eigene Collections, jedoch kann \emph{Variante A} in einem einzigen Eintrag gespeichert und geladen werden. An sich besitzt der einzelne Datensatz zwar mehr Daten, allerdings hat dies keine Auswirkung auf die Gesamtperformance. Daher würde bei einem Umbau der Aggregate für die meisten Anwendungsfälle eine einzelne Suchanfrage in N+1 Suchanfragen abgewandelt werden. Dies kann bemerkbare Auswirkungen auf die Antwortzeit unserer Applikation haben.

Im ersten Absatz wurde beschrieben, dass eine Technologie keine Auswirkung auf die Architektur haben sollte, konträr dazu ist bei der Betrachtung der Applikationsperformance dies eventuell sinnvoll. Sollte eine relationale Datenbank verwendet werden, hält sich die Performance Einbußen in Grenzen. Hingegen bei einer dokumentbasierten Datenbank ist dieser Effekt vervielfacht. Daher kann die Entscheidung über den verwendeten Aggregationsschnitt durchaus von einer konkreten Technologie beeinflusst werden. 

%TODO: ,da erstezten durch andere Wörte in der ganzen Bachelorarbeit

\subsection{Parallele Bearbeitung eines großen Aggregates}
% TODO: Datenbankanomalien erklären etc.
%TODO: Lockin method
Aufgrund dessen, dass ein Aggregat eine transaktionale Grenze abbildet, wirkt sich der Aggregationsschnitt auf die mögliche Parallelität der Anfragen aus. Ausgehend von \emph{Variante A} resultiert eine Bearbeitung der Daten des Baskets in der Speicherung des kompletten Warenkorbs in der Datenbank. Bei Schreibprozessen können Anomalien auftreten wodurch die Transaktion bzw. der Schreibvorgang abgebrochen werden muss. Eine mögliche Anomalie ist das sogenannte 'Lost Update'-Problem. Es kann auftreten wenn zwei Transaktionen den gleichen Datensatz zeitgleich bearbeiten. Sie beginnen ihre Operationen auf den selben Startzustand und schreien ihre Ergebnisse zurück in die Datenbank. Dabei kann die Änderungen einer Transaktion sofort durch eine andere überschrieben werden und dessen Datenänderungen gehen verloren.

Konkretisiert kann dieses Problem anhand des Baskets. Zwei Personen fügen beispielsweise einen neuen Artikel zu ein und demselben Warenkorb hinzu. Dies kann zum Beispiel bei einer Wishlist auftreten, welche von einer Personengruppe bearbeitet werden kann. Der erste Kunde legt hierbei ein neues Handy in den Warenkorb, wohingegen zeitgleich ein anderer Nutzer einen Fernseher hinzufügt. Beide Transaktionen starten mit einem leeren Warenkorb und schreiben in die Datenbank einen aktualisierten Datensatz mit nur einem Artikel. Daher geht die Änderung eines der zwei Personen verloren. 

Eine Lösung dieses Problems kann durch optimistischen oder pessimistischen Verfahren erzielt werden. Entweder wird im optimistischen Ansatz anhand eines Versionsfeldes überprüft, ob der Datensatz in der Zwischenzeit geändert worden ist und im Falle dessen zurückgewiesen, oder der Datensatz wird durch die pessimistischen Methode gesperrt und die zweite Transaktion muss warten bis die erste den Schreibprozess abschließt. Dementsprechend ist die zeitgleiche Bearbeitung eines Aggregats eingeschränkt. Existieren Anwendungsfälle in denen eine hohe Parallelität notwendig ist, sollten die Aggregates sofern realisierbar voneinander getrennt werden. 

\subsection{Bewertung des großen Aggregationsschnitts}

Abschließend kann anhand der vorhergehenden Überlegungen der Aggregationsschnitt von \emph{Variante A} bewertet werden.

\begin{itemize}[noitemsep,nolistsep,topsep=-2pt]
	\item \textbf{Komplexität: } {Die Handhabung und Anwendung der Businessanforderungen können maximal mühelos umgesetzt werden. Der Source Code muss durch keine komplizierteren Vorgänge ergänzt werden. }
	\item \textbf{Performance: } {Jede Anfrage benötigt das Auslesen des ganzen Baskets. Jedoch ist dies in den meisten Anwendungsfällen ohnehin von Nöten. Ein kleinere Aggregationsschnitts kann die Performance sogar verschlechtern. }
	\item \textbf{Parallelität: } {Eine zeitgleiche Bearbeitung des Baskets oder eines Objektes innerhalb ist nicht möglich.}
	\item \textbf{Anwendbarkeit: } {In Hinsicht auf die wichtigsten Anwendungsfälle erfährt der Client keine Einschränkungen und alle Businessanforderungen können erfüllt werden.}
\end{itemize}

\section{Trennung der Zahlungsinformationen von dem Basket-Aggregate}

Damit ein neues Aggregat aus der großen Basket-Klasse herausgeschnitten werden kann, wird ein Root Aggregate benötigt. Aus diesem Grund ist der erste Ansatz eine Entity vom Basket abzutrennen. Hierbei ist der PaymentProcess nur schwach an den eigentlichen Basket gebunden und mögliche Anwendungsfälle beziehen sich meist alleinig auf entweder den Warenkorb und seine Items oder den Bezahlprozess. Daher wird in der \emph{Variante C} die Aggregate in Basket und PaymentProcess, wie in Grafik \ref{fig:VarC}, aufgeteilt. 

\begin{figure}[htbp]
	\centering
	\includesvg[inkscapelatex=false, width=0.80\textwidth]{svg/VarC.svg}
	\caption{Aggregationsschnitt der Variante C}
	\label{fig:VarC}
\end{figure}

Da die Aggregates unabhängig voneinander agieren, müssen sie auch getrennt geladen und abgespeichert werden können. Folglich wird ein neues Repository benötigt, welches diese Operationen für das neue Aggregate absolviert. Die Referenz auf den PaymentProcess wird aus dem Basket entfernt und alle Funktionen, welche Aufrufe auf dieses Objekt benötigt haben, müssen entweder in die PaymentProcess-Klasse verlagert oder durch einen Domainservice durchgeführt werden. Es entstehen Auswirkungen auf die bisherige Funktionsweise und den Datenfluss der Applikation, welche in den kommenden Unterkapiteln anhand der Initiierung des Bezahlvorgangs genauer untersucht werden.

Kurze Zusammenfassung des Anwendungsfalles: 
Bevor ein Bezahlprozess gestartet werden darf, muss eine Evaluierung des Warenkorbzustands stattfinden, wie die Überprüfung der hinzugefügten Zahlungsarten, Kundendaten und errechneten Geldbeträge. Nach erfolgreicher Validierung wird der Warenkorbzustand auf 'freeze' abgeändert und der PaymentProcess kann eingeleitet werden.

\subsection{Eventuelle Konsistenz zwischen Aggregates}
%TODO: Message Broker und Listener in gls

Bei der Implementierung von Variante A ist es möglich as Einfrieren des Warenkorbs und die Initiierung innerhalb einer Transaktion erfolgt. Aufgrund der Richtlinie, dass jede Transaktion nur ein Aggregate bearbeiten darf, müssen somit diese Funktionalitäten aufgespalten werden. Das konkrete Problem, welches sich dabei ergibt, ist im Codebeispiel \ref{lst:eventualConsistency} veranschaulicht.

\begin{minipage}{\linewidth} % No pagebreak inside a minipage
	\begin{lstlisting}[caption={Getrennte Transaktionen für die Initiierung des Bezahlvorgangs}, label={lst:eventualConsistency}, language=Kotlin]
		fun initializePaymentProcess(basket: Basket, paymentProcess: PaymentProcess) {
			basket.validate()
			basket.freeze()
			basketRepository.save(basket)
			// Mögliches Interleaving anderer Operation, welche zur Abänderung des
			// Baskets und des Validierungsergebnisses führt.
			paymentProcess.initialize()
			paymentProcessRepository.save(paymentProcess)
		}
	\end{lstlisting}
\end{minipage}

%TODO: Transactional Annotation in POC Code
%TODO: Thread
Zwischen dem Persistieren des eingefrorenen Warenkorb und der Initiierung des PaymentProcesses können aufgrund von Parallelität der Warenkorb durch separate Anfragen weiterhin bearbeitet werden. Konkret könnte eine Bedingung für die Validierung sein, dass die Kundendaten nicht leer sein dürfen. Jedoch wird zwischen Zeile 4 und Zeile 7 in einem parallel laufenden Thread durch einen externen API-Aufruf die Kundendaten entfernt. Dadurch wäre die eigentliche Validierung in Zeile 2 fehlgeschlagen und die Initiierung darf nicht durchgeführt werden, jedoch wird diese dennoch in unseren Beispiel angestoßen, da davon ausgegangen wird, dass sich der Basket in der Zwischenzeit nicht verändert hat. Die diese beiden Operationen nicht als eine atomare Einheit umgesetzt ist, entsteht immer eine Möglichkeit des Interleavings zwischen den einzelnen Codezeilen.

Verdeutlicht wird dieser Effekt bei der Benutzung eines Message Brokers, falls exemplarisch die einzelnen Prozesse in verschiedenen Microservices ablaufen. Hierbei wird die konkrete Initiierung in einem anderen Service durchgeführt, welcher bei Empfang eines bestimmten Events ausgeführt wird. Folglich feuert die Funktion nach Einfrieren des Warenkorbs ein solches Event ab. Ein möglicher Programmablauf ist in Diagramm \ref{fig:VarC-Sequence} zu sehen. 

\begin{figure}[htbp]
	\centering
	\includesvg[inkscapelatex=false, width=0.80\textwidth]{svg/VarC Sequenz.svg}
	\caption{Vereinfachtes Sequenzdiagramm zur Initiierung des Bezahlvorgangs in Variante C}
	\label{fig:VarC-Sequence}
\end{figure}

Der rot markierte Bereich stellt die Zeitspanne dar, in welcher andere Prozesse den Warenkorb weiterhin bearbeiten können und dadurch die eigentliche Initiierung invalide, jedoch weiterhin durchgeführt wird. Als Folge dessen entstehen weitere Herausforderungen für das Domainmodell und die Interleavingsmöglichkeiten müssen stets in Betracht gezogen werden. Diese Art von Konsistenz wird als 'Eventuelle Konsistenz' betitelt, da es einen Zeitpunkt gibt in dem die Businessanforderungen temporär nur \emph{eventuell} erfüllt sind. Generell kann in vielen Anwendungsfällen eine kurzzeitige Abweichung der Voraussetzungen akzeptiert werden, zum Beispiel in einem Gruppenchat hat ein kurzer verzögerter Empfang der Nachrichten unter den Teilnehmern keinen großen Einfluss auf die Nutzererfahrung oder korrekte Funktionsweise der Applikation. Hingegen gilt bei der Initiierung des Zahlungsprozesses ein hoher Fokus auf fiskalisch korrekte Abarbeitung des Prozesses. Clients, welche den aktuellen Stand des Bestellprozess abfragen, erhalten bis zur Speicherung des Payment Aggregates einen veralteten Stand. Dies stellt einen bedeutenden Bruch der Invarianten dar. Folglich kann das Zusammenspiel zwischen Basket und PaymentProcess nicht durch eventuelle Konsistenz gelöst werden ohne ein hohes Risiko einzugehen.

\subsection{Atomare Transaktionen über mehrere Aggregates}

%TODO: Race Conditions in glossary
%TODO: Sehr oft "Aggregate innerhalbt einer Transaktion"
%https://softwareengineering.stackexchange.com/questions/356106/ddd-why-is-it-a-bad-practice-to-update-multiple-aggregate-roots-per-transaction   QUELLE FÜR HOST AUSSAGE
Ein weiterer Ansatz zum Bewältigen des Problems ist der Einsatz einer Transaktion über mehrere Aggregates hinweg, obwohl dies die vorher definierte Richtlinie bricht. Argumentativ muss hierzu zuerst die Frage beantwortet werden, aus welchem Grund überhaupt eine Transaktion nicht mehrere Aggregates bearbeiten darf. Durch Anpassen der Fragestellung wird dieser Aspekt klarer. 

Der Hintergedanke des Einsatzes von Aggregates ist eine Gruppierung von Klassen, welche vor und nach einer Transaktion stets konsistent sein müssen. Dies schützt die Applikation vor invaliden Zuständen und erlaubt die Annahme, dass alle Businessvoraussetzungen erfüllt sind. Kann diese Eigenschaft nicht garantiert werden, muss die Software und ihre Clients eventueller Konsistenz handhaben können. Die Aggregationsgrenzen sind somit mit der transaktionalen Konsistenz äquivalent. Dadurch entspricht die Speicherung zweier Aggregates innerhalb einer einzelnen Transaktion eine Überschreitung dieser Grenzen und lässt vermuten, dass der Aggregationsschnitt neu eingeteilt werden muss. Die Richtlinie stellt folglich lediglich ein Indiz für korrektes bzw. inkorrektes Design der Aggregates dar. Bedenklich ist alleinig, wenn die zugehörigen Aggregates auf unterschiedlichen Datenbankhosts liegen. Eine atomare Transaktion ist dadurch unmöglich, wodurch eine solche Architektur die Probleme einer eventuell konsistenten Applikation entwickelt und Interleavings auftreten können. Bei einer Microservice Architektur tritt dieser Fall oft ein, da jede einzelne Applikation meist eine eigene Datenbank auf unterschiedlichen Hosts besitzen.

%TODO: https://docs.mongodb.com/manual/core/write-operations-atomicity/
Eventuelle Konsistenz ist aus den Gründen des vorgehenden Unterkapitels nicht anwendbar und für die Checkout-Software existiert aktuell kein Anlass zur Annahme, dass zukünftig mehrere Datenbankhosts benötigt werden. Dadurch kann die Überlegung entstehen, das Einfrieren des Warenkorbs und die Initiierung des Bezahlprozesses innerhalb einer Transaktion auszuführen. Dazu ist erforderlich, dass die darunterliegende Datenbank eine Transaktion über mehrere Tabellen beziehungsweise Einträge erlaubt. In diesem Projekt wird eine MongoDB zur Datenspeicherung verwendet, welche seit Version 4 atomare Operationen auf mehrere Dokumente und Collections unterstützt. Collections sind hierbei grundlegend gleichbedeutend mit Tabellen aus relationalen Datenbanken. Unsere Applikation erfüllt dementsprechend diese Voraussetzung und eine Implementierung dieses Lösungsansatzes ist technisch möglich. Die Funktion aus dem Beispiel \ref{fig:VarC-Sequence} wird ergänzt, indem ein optimistisches Sperrverfahren auf die betroffenen Einträge in der Datenbank angewandt wird. Dadurch werden andere Zugriffe auf diesen Basket bzw. PaymentProcess erst ausgeführt, wenn die Initiierung vollständig abgeschlossen ist. Eine Möglichkeit zum Interleaving wird dadurch verhindert. Diese Vorgehensweise kann analog in ähnlichen Szenarien verwendet werden.

Weiterhin besteht die Frage, ob \emph{Variante C} ein valides Aggregationsdesign darstellt, da in vielen Funktionen eine Überschreitung der transaktionalen Grenzen stattfindet. Begründet kann diese Entscheidung dadurch, dass der Bezahlvorgang eine kritische Operation aus rechtlicher Sicht darstellt. Sobald dieser gestartet wird, muss das ganze Datenmodell stets konsistent sein. Wenn es uns nicht erlaubt sein sollte, eine Transaktion über mehrere Aggregates durchzuführen, sind somit alle Aggregationsschnitte, abgesehen von \emph{Variante A}, nicht zulässig. Eventuell kann diese Erkenntnis bereits eine vorläufige Antwort für die initiale Forschungsfrage dieser Arbeit darstellen.

Letztendlich soll die Architektur und das Datenmodell die Businessprozesse optimal unterstützen, während die Software weiterhin flexibel und performant bleibt. Die Richtlinie stellt einzig einen Leitfaden dar, um ein korrektes Design zu erleichtern, jedoch keine absolute Regel. Der Anreiz für eine genauer Unterteilung der Aggregates ist das separate Laden und zeitgleiche Bearbeiten unterschiedlicher Aggregates, welches weiterhin in \emph{Variante C} entlang der Richtlinie möglich ist. Anhand dieser Begründung wird für den Proof-of-Concept angenommen, dass ein solcher Aggregationsschnitt als Designentscheidung vertretbar ist.

%TODO: Implementierung der Variante
\subsection{Implementierung von Aggregates mit geteilten Invarianten}



\section{Verkleinerung der Aggregates durch Anpassung existierender Businessanforderungen }
\blindtext

\section{Maximale Unterteilung des Models mithilfe asynchroner Verarbeitung }
\blindtext

  
  \chapter{Implementierung des Proof-of-Concepts}

Im Verlaufe des Projekts wurde das Datenmodell mitsamt der Hexagonalen Architektur und allen relevanten Anwendungsfällen in einem Proof-of-Concept implementiert. Das Ziel dieser Software ist die Unterstreichung einer möglichen praktischen Umsetzung der Businessanforderungen unter Anwendung der gewonnen Erkenntnisse. 

Zu Beginn wurde ein neues Projekt aufgesetzt, welches benötigte Frameworks und Abhängigkeiten importiert, um \gls{Boilerplate}-Code weitestgehend zu vermeiden. Die gesamte Software \note{wurde}{Wortwiederholung} mithilfe von Kotlin entwickelt und Komponente sind analog zu einer Hexagonalen Architektur in die drei Bereiche primäre Adapter, Applikationskern und sekundäre Adapter aufgeteilt. Innerhalb des Applikationskerns befindet sich jegliche Businesslogik, sowie die notwendigen Applicationservices, Domainservices und das Domainmodell entsprechend des Domain-Driven Designs. Die umgesetzten Aggregationsschnitte belaufen sich auf die Variante A und D, sowie ihren Abwandlungen. Dieses Kapitel geht auf die Entwicklung des ersteren Ansatz mit einem einzelnen, großen Basket-Aggregate tiefer ein.

\comment{Bild der Architektur einfügen!}

\section{Design der primären Adapter}

Primäre Adapter sind Komponente, welche den Datenfluss aufgrund von einem Signal eines externen Systems initiieren. Sie stellen die grundlegenden Kommunikationsschnittstellen zwischen Clients und der Software dar. Im Proof-of-Concept fallen in diesen Bereich hauptsächlich die sogenannte Controller, welche für den Empfang von REST-API-Anfragen, der Deserialisierung übergebener Daten, sowie der Serialisierung des Antwortinhaltes zuständig sind. Zu Beginn jedes Anwendungsfalles wird ein Controller durch den Touchpoint angesprochen. Der jeweils zuständige Adapter wird aus einem Zusammenschluss von aufgerufener URL und HTTP-Methode bestimmt. 

Ein Controller beinhaltet lediglich Logik für den Empfang von Daten und der Formulierung zugehöriger Antworten. Alle Informationen, welche von außen stammen, müssen vor der Weitergabe an den Applikationskern in das Domainmodell umgewandelt werden. Ist dies nicht der Fall, hält der zentrale Teil der Software eine Abhängigkeit nach außen und das \acrlong{DIP} wird verletzt. Der Einsatz von sogenannten \acrlong{DTO}s (\acrshort{DTO}) ist gestattet, sofern diese in einem Modul des Applikationskern liegen, sodass dieser volle Kontrolle über ihre Struktur besitzt. Ein konkretes Beispiel für einen Controller ist im Codebeispiel \ref{lst:controller} abgebildet. Zur Implementierung dieser Funktionalität wurde das Framework 'Ktor' verwendet.

\begin{minipage}{\linewidth} % No pagebreak inside a minipage
\begin{lstlisting}[caption={Beispiel eines Controllers zum aktualiseren von Kundendaten}, label={lst:controller}, language=Kotlin]
put("/basket/{id}/customer") { 
	variable basketId = parseIdFromUrl("{id}")
	variable customer = request.body<Customer>()  
	variable basket = basketApiPort.setCustomer(basketId, customer)
	request.respond(HttpStatusCode.OK, basket)
}
\end{lstlisting}
\end{minipage}

\begin{itemize}[noitemsep,nolistsep]
	\item Zeile 1: Definiert die HTTP-Methode als 'PUT' und das Format der URL für diesen Endpunkt
	\item Zeile 2: Auslesen der BasketId aus der URL als Pfadparameters 
	\item Zeile 3: Deserialisierung der übertragenen Daten zu einem Customer-Objekt
	\item Zeile 4: Weitergabe der Objekte an den zuständigen Applicationservice über ein Interface
	\item Zeile 5: Antwort an die Anfrage mit HTTP-Status '200' und den geänderten Basket
\end{itemize}

Für jeden definierten Anwendungsfall ist ein entsprechender Port und Controller zuständig. Die tatsächliche Implementierung der Schnittstelle wird mittels \Gls{DI} durch das Framework 'Koin' bestimmt und geladen. Dadurch bleiben Abhängigkeiten jederzeit austauschbar und unabhängig testbar. Beispielsweise kann die korrekte Funktionsweise eines Controllers überprüft werden, indem der Applicationservice durch ein Test-Objekt ausgetauscht und Aufrufe des Objektes ausschließlich simuliert werden. Somit erfahren die einzelnen Komponenten in Testfällen keine Beeinflussung durch eventuell inkorrekt implementierten Code und Test-Fehlschläge können eindeutig einem bestimmten Abschnitt der Software zugeschrieben werden.

\section{Realisierung des Applikationskerns}

Der Applikationskern stellt das Herz der Anwendung dar. Das maßgeblich Ziel einer Hexagonalen Architektur ist es, dass Zentrum komplett von äußeren Modulen zu entkoppelt. Somit können Adapter jederzeit ausgetauscht werden, ohne hierbei die Businesslogik anpassen zu müssen. Die Kommunikation mit dem Kern geschieht ausschließlich über Interfaces, den Ports.

\subsection{Definition von Applicationservices anhand ihrer Aufgaben}

Entsprechend der Definition im Domain-Driven Design liegen Applicationservice eine logische Ebene höher als die Domainservice. Der Unterschied zwischen den beiden ist ihre Einsichten in die Domain. Ein Applicationservice darf kein direktes Domainwissen verwirklichen. Dazu zählen Invarianten oder die Umsetzung von Businessanforderungen.
Meist beinhalten diese Service folgende Anweisungen:

\begin{itemize}[noitemsep,nolistsep]
	\item Starten und schließen einer Transaktion
	\item Laden von Aggregates aus Repositories
	\item Einfache Ablaufsteuerung
	\item Aufruf von Funktionen auf den Aggregates oder Domainservices
	\item Umsetzung von technologischen Komponenten wie Event-Listener oder Caching
\end{itemize}

Ein simples Beispiel bietet uns der BasketItem-Applicationservice. Dieser wird aufgerufen, sofern Änderungen an den BasketItems durchgeführt werden sollen. Der in Figur \ref{lst:applicationservice} dargestellte Code behandelt die Entfernung eines BasketItem aus dem Warenkorb.

\begin{minipage}{\linewidth} % No pagebreak inside a minipage
	\begin{lstlisting}[caption={Eine Beispielsfunktion des BasketItem-Applikationservice}, label={lst:applicationservice}, language=Kotlin]
function removeBasketItem(BasketId basketId, BasketItemId basketItemId) {
	variable basket = basketRepository.findById(basketId)
	basket.removeBasketItem(basketItemId)
	basketStorageService.store(basket)
}
	\end{lstlisting}
\end{minipage}

\begin{itemize}[noitemsep,nolistsep]
	\item Zeile 2: Laden eines Baskets anhand seiner Id durch einen Domainservice. 
	\item Zeile 3: Aufruf des Basket-Aggregate-Root zum Entfernen des übergebenen Items. Innerhalb dieser Funktion werden zusätzlich Aufgaben erledigt, wie das Neuberechnen des Gesamtpreises. Sollte die Kalkulation zu komplex ausfallen, kann hier ein Berechnungsservice als Parameter übergeben werden. Somit ist der Basket für seine eigene Konsistenz verantwortlich. 
	\item Zeile 4: Speichern des Warenkorbs mit den abgeänderten Daten.
\end{itemize}

\subsection{Aufteilen der Businesslogik zwischen Domainservices und Datenmodell}

Aufgrund der Tatsache, dass die Software lediglich ein Aggregat und somit auch nur ein Aggregate-Root besitzt, müssen alle Operationen auf den Attributen des Baskets durch eine Methode im Basket selbst geschehen. Anwendungsfälle, die es erfordern tiefer gelegene Objekte anzupassen, werden durch eine Kette von Funktionsaufrufen umgesetzt. Hierbei können alle notwendigen Invarianten überprüft werden, da das komplette Datenmodell geladen ist. Zur Trennung der Datenhaltung von den Funktionalitäten implementiert der Warenkorb ein Interface, welches jegliche von außen notwendigen Schnittstellen beinhaltet. Dadurch kann das darunterliegende Datenmodell ausgetauscht oder in Tests simuliert werden. Als Veranschaulichung für das Abändern der Fulfillment Methode dient der Codeauszug \ref{lst:basket}.

\begin{minipage}{\linewidth} % No pagebreak inside a minipage
	\begin{lstlisting}[caption={Setzen der Fulfillment Methode im Basket Aggregate}, label={lst:basket}, language=Kotlin]
function setFulfillment(Fulfillment fulfillment, FulfillmentPort fulfillmentPort) {
	validateIfModificationIsAllowed()
	variable availableFulfillment = fulfillmentPort.getAvailableFulfillment(outletId)
	throwIf(availableFulfillment.doesNotContain(fulfillment)) {
		IllegalModificationError("$fulfillment is not avaiable")
	}
	this.fulfillment = fulfillment
}
	\end{lstlisting}
\end{minipage}

\begin{itemize}[noitemsep,nolistsep]
	\item Zeile 2: Überprüfung, ob der Basket aktuell Änderungen zulässt anhand des Status. Stellt eine Invariante des Datenmodells dar.
	\item Zeile 3: Laden aller verfügbarer Fulfillment Methoden für diesen Basket durch einen sekundären Adapter. Die Kommunikation mit dem Adapter erfolgt über einen Port.
	\item Zeile 4-6: Falls der neue Wert nicht unter den verfügbaren Fulfillment ist, wird der Aufruf zurückgewiesen und eine entsprechende Fehlermeldung an den Client durch den Controller geliefert.
	\item Zeile 7: Überschreiben des alten Wertes. Dieser Punkt wird nicht erreicht, wenn zuvor eine Businessanforderung gescheitert ist.
\end{itemize}

Aufgaben, welche nicht direkt einem Objekt zugewiesen werden können oder mehrere Aggregates betreffen sind in Domainservices zu implementieren. Diese liegen auf der gleichen konzeptionellen Ebene wie das Datenmodell und dürfen Businesslogik enthalten. Beispielsweise existiert im Proof-of-Concept ein Domainservice für die Abwicklung des Bezahlverfahren, wie in Code-Ausschnitt \ref{lst:domainservice} aufgeführt. 

\begin{minipage}{\linewidth} % No pagebreak inside a minipage
	\begin{lstlisting}[caption={Ausführung des Bezahlvorgangs in einem Domainservice}, label={lst:domainservice}, language=Kotlin]
function executePaymentProcessAndFinalizeBasket(BasketId basketId) {
	variable basket = basketRepository.findById(basketId)
	throwIf(basket.isNotFrozen() or basket.paymentIsNotInitialized()) {
		IllegalModificationError("cannot cancel payment process")
	}
	variable externalPaymentRef = basket.getExternalPaymentRef()
	paymentPort.executePayment(externalPaymentRef)
	basket.executePayments() and basket.finalize()
	basketStorageService.store(basket)
	createOrderAfterFinalization(basket)
}
	\end{lstlisting}
\end{minipage}

\begin{itemize}[noitemsep,nolistsep]
	\item Zeile 2: Laden des aktualisierten Baskets aus dem Repository.
	\item Zeile 3-5: Weist die Durchführung zurück, sofern der Basket sich nicht in den erwarteten Zustand befindet. Dies kann auftreten, wenn die REST-API aufgerufen worden ist, ohne dass ein Zahlungsprozess zuvor gestartet wurde.
	\item Zeile 6-8: Durchführung des Bezahlvorgangs. Die erforderliche Aufrufreihenfolge stellt einen Teil des Domainwissens dar und begründet die Zuteilung der Klasse in die Gruppe der Domainservices.
	\item Zeile 9: Speichern des angepassten Warenkorbs.
	\item Zeile 10: Erstellen eines Bestellvorgangs durch einen separaten Domainservice.
\end{itemize}

\section{Anbinden externer Systeme und Datenbanken durch sekundäre Adapter}

Das, in der Planungsphase erstellte, Context-Diagramm \ref{fig:ContextDiagramm} zeigt verschiedenste Systeme mit denen die Anwendung zum Erfüllen ihrer Aufgaben kommunizieren muss. Für diesen Zweck wurden zwei Gruppen von Komponente eingeführt. Die Service, welche Aufrufe von externen API-Schnittstellen simulieren, und die sekundären Adapter. Letztere implementieren ein vom Applikationskern definiertes Interface, damit eine Brücke zwischen dem Domainkern und den APIs der Nachbarsysteme gewährleistet ist. Analog zu den primären Adapter, existieren im Anwendungskern keine direkten Abhängigkeiten zu diesem Teil der Software.

Ein Spezialfall sind hierbei die Service für das Erfragen der aktuellen Preis- bzw. Artikelinformationen. Aufgrund von Performance-Verbesserungen wurde eine Abwandlung der Adapter mit Caching-Funktion zusätzlich erstellt. Der normal fungierende Service ruft den zugehörigen API-Service auf, wohingegen der Caching-Adapter zuerst den Preis bzw. Artikel aus dem Cache lädt und zurückgibt, sofern dieser noch als aktuell markiert ist. Sollte der Eintrag veraltet sein, wird der eigentliche Adapter angesprochen, um die zum jetzigen Zeitpunkt validen Daten aus dem externen System zu erfragen. Im Beispiel \ref{lst:adapter} ist die, für das Aktualisieren des Preises zuständige, Klasse  abgebildet.

\begin{minipage}{\linewidth} % No pagebreak inside a minipage
	\begin{lstlisting}[caption={Preisadapter mit Caching-Funktion}, label={lst:adapter}, language=Kotlin]
class CachedPriceAdapter {
	
	variable PriceRepository priceRepository 
	variable PriceAdapter priceAdapter
	
	function fetchPrice(PriceId priceId) returns Price {
		return priceRepository.getAndUpdateIfInvalid(priceId, fallback = {
			priceAdapter.fetchPrice(priceId)
		})
	}
}
	\end{lstlisting}
\end{minipage}

\begin{itemize}[noitemsep,nolistsep]
	\item Zeile 3-4: Der Cache-Preisadapter hat eine Abhängigkeit zum normalen Preisadapter und zu einem Repository zum Abrufen des zwischengespeicherten Preises
	\item Zeile 7: Abfragen des Preises aus dem Cache-Repository. Sollte der Preis invalide sein wird Zeile 8 ausgeführt.
	\item Zeile 8: Weiterleitung der Anfrage an den zuständigen Adapter. Das Ergebnis wird mit einem aktuellen Zeitstempel im Cache abgelegt.
\end{itemize}

Zusätzlich zu diesen Adaptern gehören ebenfalls die Repositories in diesem Bereich der Hexagonalen Architektur. Sie managen den Zugriff auf die Datenbank und alle Funktionalitäten, welche in dieses Aufgabengebiet fallen. In dem Proof-of-Concept wird das gesamte Basket-Aggregate abgespeichert. Aus diesem Grund wird nur ein Repository benötigt, welches eine grundlegende Suchfunktion mithilfe der BasketId und eine Speicherfunktion anbietet. Abgesehen davon gehören noch die zwei Caching-Repositories zu dieser Gruppe. 
  
  
%% FAZIT %% 

\chapter{Fazit und Empfehlungen}

Der Aufbau einer Hexagonalen Architektur unterstützt bei der Entkopplung des Applikationskerns und ermöglicht das instinktiv Einhalten der SOLID-Prinzipien. In dem Proof-of-Concept hat eine Einteilung in Adapter, Ports und Businesslogik den Entwicklungsprozess erleichtert. Zusätzlich lässt sich diese Softwarestruktur natürlich mit einem Domain-Driven Ansatz kombinieren. Die Hexagonale Architektur bildet ein stabiles und erweiterbares Fundament für die Checkout-Applikation.

Ein effektiver Aggregationsschnitt ist anhand seiner erhöhten Skalierbarkeit und Performance erkennbar. Damit ein Umbau der aktuellen Produktivsoftware zu empfehlen ist, sollten dementsprechend diese Qualitätsmerkmale positiv betroffen sein. 

Mithilfe der Performance-Analyse kann ein Performance-Gewinn hierbei nicht verzeichnet werden. Dies lässt sich auf die vielzähligen Datenbankoperationen bei einem aufgeteilten Aggregationsschnitt zurückführen, welche aufgrund der starke Kopplung einzelner Entitäten untereinander entstehen. In vielen Softwareprojekten existieren kaum Invarianten, die sich über das ganze Domain-Modell spannen. Oftmals ist deshalb eine Trennung der Klassen voneinander problemlos möglich und ein verbessertes Aggregationsdesign kann dadurch erreicht werden. Der Warenkorb ist allerdings ein enger Verbund aus Businessrichtlinien und seine transaktionale Konsistenz muss wegen fiskalischen Anforderungen stets vorliegen. Eine allumfassender Aggregationsschnitt ist aus Businesssicht somit ebenfalls sinnvoll.

Zudem ist die mögliche parallele Bearbeitung eines Warenkorbs zum jetzigen Zeitpunkt nicht notwendig. Weshalb aus diesem Aspekt keine negativen Einflüsse durch den großen Aggregationsschnitt entstehen. Sofern zukünftig zeitgleiche Modifikationen an einem einzelnen Basket durch verschiedene Nutzer ein gängiger Anwendungsfall darstellt, müssen die Artikel als eigenständiges Aggregate designt werden. Eine verringerte Performance und erhöhte Softwarekomplexität gehen im Gegenzug mit diesem Domain-Modell einher.

Das Abspalten der Preiskalkulation vom Warenkorbs ist denkbar, um Berechnungszeiten einzusparen. Die Komplexität des Sourcecodes steigt in diesem Fall leicht an. Sollten die Touchpoints bei der Mehrheit der API-Anfragen jedoch auch zugleich die neu kalkulierten Werte erwarten, verfallen die Performance-Verbesserungen. \note{Weitere Designvariationen sind, anhand der, in den jeweiligen Unterkapiteln besprochenen, Auswirkungen zur Verbesserung der Qualitätsmerkmalen ungeeignet.}{Komma, ja oder nein}

\textbf{Schlussendlich fallen die Argumente für einen Umgestaltung der Aggregationsaufteilung in Zusammenhang mit den aktuellen Businessanforderungen zu schwach aus, sodass ein Neudesign der Applikation nicht empfehlenswert ist.}

Allgemein kann sich in Softwareprojekten ein idealer Aggregationsschnitt mithilfe einer detaillierten Analyse der Anwendungsfälle und Integritätsgrenzen herauskristallisieren. Die ermittelten Invarianten zwischen Objekte bestimmen maßgeblich umsetzbare Designansätze. Ferner können verwendete Technologien durchaus einen Einfluss auf die Architektur der Software haben, allerdings sollte dies mit dem Bewusstsein geschehen, dass sich diese zeitnahe ändern können und weiterhin bei deren Einbindung in den Entscheidungsprozess ein Risiko entsteht. Anforderungen an die Applikation und in der Entwicklungsphase entstehende Kompromisse besitzen oftmals eine größere Priorität im Vergleich zu den theoretischen Prinzipien des Softwaredesigns. Diese bieten zwar Richtlinien für eine langlebige Anwendung, allerdings sind sie nicht immer die optimale Lösung für ein konkretes Problem. Der zugrundeliegenden Antrieb für ihre Einhaltung sollte hinterfragt und bewusst gemacht werden. Nichtsdestotrotz ist in den meisten Fällen ein kleinerer Aggregationsschnitt zu bevorzugen, da viele Applikationen kaum starke Invarianten besitzen. Eine zukunftssicheres Software gelingt somit unter Einhaltung der gängigen Richtlinien eines Domain-Driven Designs in Kombination mit etablierten Designprinzipien und intuitiven Lösungsansätzen von erfahrenen Entwickler:innen.  

\comment{Zu wenig. Mehr ausformulieren. Mehr Zukunftsausblicke etc. Abruptes Ende...}
  
  \backmatter

	\def\insideAnhang{0}

\newenvironment{anhang}{%
	\ifnum\insideAnhang=1%
	\errhelp={Let other blocks end at the beginning of the next block.}
	\errmessage{Nested Alpha section not allowed}
	\fi%
	\def\insideAnhang{1}
	
	\renewcommand\figurename{%
		\ifnum\insideAnhang=1% 
		Anhang\else%
		Abbildung\fi%
	}%
	
	\ifnum\insideAnhang=1%
	\def\table{\def\figurename{Anhang}\figure}
	\let\endtable\endfigure
	\fi%
	
	\renewcommand\thefigure{%
		\ifnum\insideAnhang=1% 
		\Alph{figure}\else%
		\arabic{figure}\fi%
	}%
	
}{%
	\def\insideAnhang{0}
}%

%%%%%%%%%%%%%%%%%%%%%%%%%%%%%%%%%%%%%%%%%%%%%%%%
%%%%%%%%%%%%%%%%%%%%%%%%%%%%%%%%%%%%%%%%%%%%%%%%
%%%%%%%%%%%%%%%%%%%%%%%%%%%%%%%%%%%%%%%%%%%%%%%%
%%%%%%%%%%%%%%%%%%%%%%%%%%%%%%%%%%%%%%%%%%%%%%%%

\begin{anhang}

\chapter{Anhang}

\section{Sourcecode des Proof-of-Concepts}
\label{label:sourcecode}

Die Binärdateien des Projekts wurden zur Versionierung in ein Git-Repository hinterlegt. Diese umfassen die Definition der Lasttests, API-Beschreibung, und den Quelltext, inklusive der analysierten Aggregationsschnitte. Das Repository kann unter dem Link '\url{https://github.com/Thalmaier/bachelor-thesis-checkout-poc}' aufgerufen werden. Alternativ ist in Anhang \ref{fig:Github} ein QR-Code abgebildet.

\begin{figure}[h]
	\centering
	\vspace{0.8cm}
	\includesvg[inkscapelatex=false, width=100px, height=100px]{svg/github qr code.svg}
	\caption{QR-Code des GitHub Repositories}
	\small URL: \hspace{0.3mm} \url{https://github.com/Thalmaier/bachelor-thesis-checkout-poc}
	\label{fig:Github}
\end{figure}

\newpage

\section{Aktivitätsdiagramme der Anwendungsfälle}

\begin{figure}[h!]
	\centering
	\includesvg[inkscapelatex=false, width=\textwidth]{svg/AD Basketcreation.svg}
	\caption{Aktivitätsdiagramm für die Erstellung eines Baskets}
	\label{fig:SL-Basketcreation}
\end{figure}

\begin{figure}[h!]
	\centering
	\includesvg[inkscapelatex=false, width=\textwidth]{svg/AD Basketstornierung.svg}
	\caption{Aktivitätsdiagramm für die Stornierung eines Baskets}
	\label{fig:SL-Basketstornierung}
\end{figure}

\begin{figure}[h!]
	\centering
	\includesvg[inkscapelatex=false, width=\textwidth]{svg/AD Checkoutdata.svg}
	\caption{Aktivitätsdiagramm für das Setzen der Checkout Daten}
	\label{fig:SL-Checkoutdata}
\end{figure}

\begin{figure}[h!]
	\centering
	\includesvg[inkscapelatex=false, width=\textwidth]{svg/AD PutBezahlmethode.svg}
	\caption{Aktivitätsdiagramm für das Hinzufügen einer Bezahlmethode}
	\label{fig:SL-PutBezahlmethode}
\end{figure}

\phantom{}
\newpage

\section{API-Endpunkte}

Die OpenAPI Definition der API ist in Anhang \ref{fig:Github} hinterlegt. Zudem bietet bietet Anhang \ref{fig:REST-API} eine grafische Übersicht der Endpunkte für Variante A des Aggregate-Designs.

\begin{figure}[h!]
	\centering
	\includesvg[inkscapelatex=false, width=\textwidth, height=0.85\textheight]{svg/REST API.svg}
	\caption{REST-API der Checkout-Software für diesen Proof-of-Concept}
	\label{fig:REST-API}
\end{figure}

\newpage
\section{Vollständiges Datenmodell des Proof-of-Concepts} \label{label:Daten-Modell}

\groupedDomainModell{Basket}{
	\item \textbf{BasketId: } {Eindeutige Identifikation des Baskets zur Referenzierung durch die Touchpoints.}
	\item \textbf{OutletId: } {Eine Referenz zugehörig zu dem Markt oder Onlineshop, durch welchen der Basket angelegt wurde. Unerlässlich für die Bestimmung von unter anderem Lagerbeständen, Lieferzeiten, Fulfillment-Optionen und Versandkosten.}
	\item \textbf{BasketStatus: } {Repräsentiert den aktuellen Zustands des Baskets. Mögliche Werte sind 'OPEN', 'FROZEN', 'FINALIZED' und 'CANCELED'.} 
	\item \textbf{Customer: } {Speichert Kundendaten (IdentifiedCustomer) oder Session-Informationen (SessionCustomer).} 
	\item \textbf{FulfillmentType: } {Lieferart, wie 'PICKUP' oder 'DELIVERY'.} 
	\item \textbf{BillingAddress: } {Adresse für die Rechnungserstellung.} 
	\item \textbf{ShippingAddress: } {Adresse für die Warenlieferung.} 
	\item \textbf{BasketItems: } {Liste aller enthaltenen Produkten und ihren zugehörigen Informationen.}
	\item \textbf{BasketCalculationResult: } {Beinhaltet die berechneten Werte des Basket, wie Nettobetrag, Bruttobetrag und Mehrwertsteuer. Die Speicherung dieser Werte wäre technisch nicht notwendig, spart aber an Rechenzeit ein, da nicht bei jeder Abfrage des Basket dieser Wert neu berechnet werden muss.}
	\item \textbf{PaymentProcess: } {Bindet alle Informationen zur erfolgreichen Abwicklung des Zahlungsprozesses.}
	\item \textbf{Order: } {Speichert eine Referenz auf die Bestellung für einen Basket. Wird erst nach Zahlungsabschluss befüllt.}  
}

\groupedDomainModell{SessionCustomer}{
	\item \textbf{SessionID: } {Eindeutige ID zur Zuweisung einer Session im Onlineshop zum zugehörigen Basket. Notwendig, um eine Einkaufmöglichkeit für anonyme Kunden zu bieten.}
}

\groupedDomainModell{IdentifiedCustomer}{
	\item \textbf{Name: } {Enthält den Vor- und Nachnamen als eigenes Datenkonstrukt.}
	\subDomainModell{
		\item \textit{FirstName: } {Vorname des Kunden.}
		\item \textit{LastName: } {Nachname des Kunden.}
	}
	\item \textbf{E-Mail: } {E-Mail des Kunden.}
	\item \textbf{CustomerTaxId: } {Die Steuernummer des Kunden. Relevant aus Sicht der Rechnungsabwicklung und für den Ausdruck der Rechnung.} 
	\item \textbf{BusinessType: } {Bestimmt ob Kunde als Business-to-Customer (B2C) oder Business-to-Business (B2B) gilt.}
	\item \textbf{CompanyName: } {Firmenname des Kunden. Kann optional angegeben werden oder ist verpflichtend für einen B2B-Kunden.}
	\item \textbf{CompanyTaxId: } {Steuernummer der Firma eines B2B-Kunden.}
}

\groupedDomainModell{BasketItem}{
	\item \textbf{Id: } {Eindeutige Referenz auf den Warenkorbeintrag.}
	\item \textbf{Product: } {Beinhaltet alle Produktinformationen, welche durch die Touchpoints benötigt werden.}
	\item \textbf{Price: } {Aktueller Preis des zugehörigen Products. Kann sich zeitlich ändern, muss daher durch eine Businessfunktion aktualisiert werden. }
	\item \textbf{ShippingCost: } {Betrag der Lieferkosten des Items.}
	\item \textbf{BasketItemCalculationResult: } {Speichert die Bruttokosten des Produktes, die errechneten Nettokosten, Lieferkosten und den Gesamtpreis.}
}

\groupedDomainModell{Product}{
	\item \textbf{Id: } {Eindeutige Referenz des Products im externen System.}
	\item \textbf{Name: } {Textuelle Produktbezeichnung des Products.}
	\item \textbf{Vat: } {Mehrwertsteuerinformationen des Products.}
	\item \textbf{UpdatedAt: } {Zeitstempel notwendig für die Aktualisierungsfunktion der Artikelinformationen.}
}

\groupedDomainModell{Vat}{
	\item \textbf{Sign: } {Identifizierung des Steuertyps, abhängig von jeweiligen Prozentsatz und Land.}
	\item \textbf{Rate: } {Prozentualer Wert der Mehrwertsteuer, wie beispielsweise '19\%'.}
}

\groupedDomainModell{Price}{
	\item \textbf{PriceId: } {Setzt sich zusammen aus der ProductId und der OutletId.}
	\item \textbf{GrossAmount: } {Bruttobetrag mit Währung.}
	\item \textbf{UpdatedAt: } {Zeitstempel notwendig für die Aktualisierungsfunktion des Preises.}
}

\groupedDomainModell{BasketItemCalculationResult}{
	\item \textbf{ItemCost: } {Beinhaltet Netto, Brutto und VAT Informationen in Form eines CalculationResults.}
	\item \textbf{ShippingCost: } {Betrag der Lieferkosten.}
	\item \textbf{TotalCost: } {Zusammengerechnete Werte der einzelnen Preise im Form eines CalculationResults.}
}

\groupedDomainModell{CalculationResult}{
	\item \textbf{GrossAmount: } {Bruttobetrag mit Währung.}
	\item \textbf{NetAmount: } {Nettobetrag mit Währung.}
	\item \textbf{VatAmounts: } {Eine zusammengebautes Set aus VatAmounts der Preise der BasketItems. Benötigt, da Vats mit unterschiedlichen Prozentbeträgen rechtlich nicht kombiniert werden dürfen.}
}

\groupedDomainModell{VatAmount}{
	\item \textbf{Sign: } {Identifizierung des Steuertyps, abhängig von genauen Prozentsatz und zugehörigen Land.}
	\item \textbf{Rate: } {Prozentualer Wert der Mehrwertsteuer.}
	\item \textbf{Amount: } {Berechneter Betrag der Mehrwertsteuer zugehörig zu einem Bruttobetrag.}
}

\groupedDomainModell{BasketCalculationResult}{
	\item \textbf{GrandTotal: } {Betrag der finalen Gesamtkosten des ganzen Baskets.}
	\item \textbf{NetTotal: } {Fasst alle Nettobeträge zusammen in einem einzelnen Betrag.}
	\item \textbf{ShippingTotal: } {Fasst alle Lieferkosten zusammen in einem einzelnen Betrag.}
	\item \textbf{VatAmount: } {Rechnet alle Vats zusammen, welche das gleiche Sign besitzen.}
}

\groupedDomainModell{PaymentProcess}{
	\item \textbf{BasketId: } {Id des zugehörigen Baskets.}
	\item \textbf{ExternalPaymentRef: } {Referenz auf den Bezahlvorgangs im externen System. Anfangs leer bis zur Initiierung des Payments.}
	\item \textbf{AmountToPay: } {Betrag der insgesamt bezahlt werden muss. Entspricht dem GrandTotal des Baskets.}
	\item \textbf{AmountPayed: } {Rechnet alle Payments zusammen und bestimmt in welchem Maße der Basket bereits bezahlt ist.}
	\item \textbf{AmountToReturn: } {Falls der bezahlte Betrag größer ist als gefordert, wird dieser Wert berechnet. Repräsentiert den Betrag, welcher durch das System zurückgegeben werden muss.}
	\item \textbf{PaymentProcessStatus: } {Status wieweit der der AmountToPay bezahlt ist. Kann die Werte 'TO\_PAY', 'PARTIALLY\_PAID' und 'PAID' annehmen.}
	\item \textbf{Payment: } {Liste aller Payments zugehörig zu diesem Prozess.}
}

\groupedDomainModell{Payment}{
	\item \textbf{PaymentId: } {Die Id der Zahlung.}
	\item \textbf{PaymentMethod: } {Bezahlungsart, wie Gutschein oder Barzahlung.}
	\item \textbf{PaymentStatus: } {Aktueller Zustand des Payments. Mögliche Werte entsprechen 'SELECTED', 'INITIALIZED', 'EXECUTED', 'CANCELED'. Ein Payment ist bei Hinzufügung im Status 'SELECTED'.}
	\item \textbf{AmountSelected: } {Betrag, welcher durch dieses Payment bezahlt werden soll. Falls dieser Wert leer ist, wird der gesamte Warenkorb durch dieses Payment bezahlt.}
	\item \textbf{AmountUsed: } {Betrag wie viel insgesamt durch dieses Payment abgedeckt wurde, falls nur ein Bruchteil des AmountSelectes benötigt wird.}
	\item \textbf{AmountOverpaid: } {Berechnet durch Subtraktion von AmountSelected und AmountUsed.}
}

\groupedDomainModell{Order}{
	\item \textbf{OrderRef: } {Referenz auf die Bestellung des Warenkorbs. Wird bei Abschluss des Zahlungsprozesses gesetzt.}
}

\clearpage
\section{Klassendiagramme des Datenmodells}

\vspace{1cm}
\begin{figure}[htbp]
	\centering
	\includesvg[inkscapelatex=false, width=0.95\textwidth]{svg/ValueObjectBasketDiagram.svg}
	\caption{Klassendiagramm eines Baskets}
	\label{fig:VO-Basket}
\end{figure}

\begin{figure}[htbp]
	\centering
	\includesvg[inkscapelatex=false, width=0.95\textwidth]{svg/ValueObjectCustomerDiagram.svg}
	\caption{Klassendiagramm des Customer Value Objects}
	\label{fig:VO-Customer}
\end{figure}

\begin{figure}[htbp]
	\centering
	\includesvg[inkscapelatex=false, width=0.95\textwidth]{svg/ValueObjectPaymentDiagram.svg}
	\caption{Klassendiagramm des Payment Process}
	\label{fig:VO-Payment}
\end{figure}

\section{Ergebnisse des Lasttests} \label{label:Lasttests}

Der Lasttest wurde mithilfe der Software 'JMeter' durchgeführt. Die Vorlage der jeweiligen Testfälle sind im Repository des Proof-of-Concepts unter Anhang \ref{fig:Github} zu finden. Ein Durchlauf bezieht sich auf einen typischen User Story, welche folgende Aspekte beinhaltet: Erstellen eines Baskets, dreimaliges Hinzufügen von Artikeln, Setzen der Checkout-Daten, zweimaliges Abrufen des Warenkorbs, Hinzufügen eines Payments und das Initiieren inklusive Durchführen des Bezahlvorgangs. 

Folgende Abkürzungen wurde zur Übersichtlichkeit genutzt:
\begin{itemize}[noitemsep,nolistsep]
	\item \textbf{A}: 'Variante A' des Aggregationsschnittes
	\item \textbf{D}: 'Variante D' des Aggregationsschnittes
	\item \textbf{M}: Verwendung des Datenbankmanagementsystems MongoDB
	\item \textbf{P}: Verwendung des Datenbankmanagementsystems PostgreSQL
	\item \textbf{F}: Kurz für 'Flag'. Die Kalkulation des Gesamtpreises geschieht erst bei expliziter Abfrage
	\item \textbf{C}: Die Kalkulation des Gesamtpreises findet umgehend bei Anpassungen von relevanten Werten statt.
	\item \textbf{AZ}: Kurz für 'Ablaufzeit'. Angabe in Millisekunden. Die Zeit für einen einzelnen Ablauf des Anwendungsfalles. 
\end{itemize}


\begin{landscape}
\begin{table}[h!]
	\centering
	\small
	\vspace{1cm}
	\begin{tabular}{ |c|c|c|c|c|c|c|c|} 
		\hline
		Name & Anzahl Durchläufe & min. AZ & max. AZ & durchschn. AZ & Median der AZ & Durchläufe pro Sekunde & Testdauer in Millisekunden \\ 
		\hline
		Variante A-M & 100 & 16 & 53 & 28,66 & 29 & 92,08 & 1086,00 \\
		\hline
		Variante A-M & 100 & 17 & 51 & 24,70 & 25 & 88,42 & 1131,00 \\
		\hline
		Variante A-M & 100 & 16 & 53 & 28,19 & 25 & 91,07 & 1098,00 \\
		\hline
		Variante A-M & 1000 & 16 & 83 & 43,56 & 42 & 197,36 & 5067,00 \\
		\hline
		Variante A-M & 1000 & 15 & 92 & 41,56 & 40 & 207,13 & 4828,00 \\
		\hline
		Variante A-M & 1000 & 16 & 86 & 41,53 & 41 & 205,00 & 4878,05 \\
		\hline
		Variante A-M & 10000 & 15 & 78 & 37,67 & 36 & 259,24 & 38574,00 \\
		\hline
		Variante A-M & 10000 & 15 & 76 & 36,86 & 35 & 265,32 & 37690,00 \\
		\hline
		Variante A-M & 10000 & 15 & 78 & 36,78 & 35 & 266,16 & 37571,00 \\
		\hline
		Variante D-FM & 100 & 21 & 55 & 29,10 & 29 & 88,18 & 1134,00 \\
		\hline
		Variante D-FM & 100 & 22 & 58 & 31,03 & 30 & 86,43 & 1157,00 \\
		\hline
		Variante D-FM & 100 & 22 & 60 & 33,25 & 31 & 86,73 & 1153,00 \\
		\hline
		Variante D-FM & 1000 & 21 & 75 & 42,46 & 41 & 199,60 & 5010,02 \\
		\hline
		Variante D-FM & 1000 & 22 & 83 & 44,97 & 44 & 189,47 & 5278,00 \\
		\hline
		Variante D-FM & 1000 & 21 & 80 & 43,41 & 42 & 195,69 & 5110,00 \\
		\hline
		Variante D-FM & 10000 & 21 & 73 & 37,00 & 35 & 264,63 & 37788,00 \\
		\hline
		Variante D-FM & 10000 & 20 & 97 & 36,82 & 35 & 266,36 & 37543,01 \\
		\hline
		Variante D-FM & 10000 & 20 & 72 & 36,83 & 35 & 265,82 & 37619,99 \\
		\hline
		Variante D-CM & 100 & 19 & 54 & 25,22 & 24 & 90,33 & 1107,00 \\
		\hline
		Variante D-CM & 100 & 18 & 60 & 28,40 & 27 & 90,42 & 1106,00 \\
		\hline
		Variante D-CM & 100 & 18 & 63 & 27,24 & 28 & 91,32 & 1095,00 \\
		\hline
		Variante D-CM & 1000 & 19 & 72 & 39,67 & 38 & 211,77 & 4722,00 \\
		\hline
		Variante D-CM & 1000 & 18 & 69 & 40,13 & 39 & 211,46 & 4729,00 \\
		\hline
		Variante D-CM & 1000 & 17 & 79 & 38,92 & 38 & 215,52 & 4640,00 \\
		\hline
		Variante D-CM & 10000 & 17 & 93 & 32,82 & 31 & 297,78 & 33582,00 \\
		\hline
		Variante D-CM & 10000 & 17 & 89 & 33,17 & 31 & 293,90 & 34025,00 \\
		\hline
		Variante D-CM & 10000 & 17 & 91 & 33,42 & 31 & 292,53 & 34185,00 \\
		\hline
	\end{tabular}
	\caption{Analyseergebnis des Lasttests der verschiedenen Variationen in Kombination mit MongoDB}
	\label{fig:performance-mongo}
\end{table}
\end{landscape}

\begin{landscape}
	\begin{table}[h!]
		\centering
		\small
		\vspace{1cm}
		\begin{tabular}{ |c|c|c|c|c|c|c|c|} 
			\hline
			Name & Anzahl Durchläufe & min. AZ & max. AZ & durchschn. AZ & Median der AZ & Durchläufe pro Sekunde & Testdauer in Millisekunden \\ 
			\hline
			Variante A-P & 100 & 37 & 65 & 48,08 & 48 & 76,75 & 1303,00 \\
			\hline
			Variante A-P & 100 & 34 & 77 & 55,08 & 55 & 73,58 & 1359,00 \\
			\hline
			Variante A-P & 100 & 35 & 80 & 55,04 & 54 & 70,95 & 1409,45 \\
			\hline
			Variante A-P & 1000 & 32 & 149 & 60,15 & 57 & 148,39 & 6739,00 \\
			\hline
			Variante A-P & 1000 & 33 & 133 & 60,02 & 57 & 147,21 & 6793,00 \\
			\hline
			Variante A-P & 1000 & 34 & 118 & 59,83 & 59 & 148,35 & 6741,00 \\
			\hline
			Variante A-P & 10000 & 33 & 93 & 57,28 & 57 & 170,35 & 58703,02 \\
			\hline
			Variante A-P & 10000 & 33 & 109 & 55,62 & 55 & 175,18 & 57083,00 \\
			\hline
			Variante A-P & 10000 & 33 & 138 & 55,61 & 55 & 174,23 & 57395,99 \\
			\hline
			Variante D-FP & 100 & 42 & 93 & 63,38 & 63 & 69,93 & 1430,00 \\
			\hline
			Variante D-FP & 100 & 42 & 90 & 63,53 & 63 & 70,13 & 1426,00 \\
			\hline
			Variante D-FP & 100 & 39 & 169 & 69,44 & 64 & 69,49 & 1439,00 \\
			\hline
			Variante D-FP & 1000 & 39 & 170 & 62,82 & 60 & 143,66 & 6961,00 \\
			\hline
			Variante D-FP & 1000 & 37 & 148 & 60,63 & 59 & 148,13 & 6751,00 \\
			\hline
			Variante D-FP & 1000 & 38 & 88 & 61,54 & 60 & 143,55 & 6966,00 \\
			\hline
			Variante D-FP & 10000 & 37 & 135 & 56,06 & 55 & 174,40 & 57338,00 \\
			\hline
			Variante D-FP & 10000 & 41 & 113 & 58,71 & 58 & 166,68 & 59994,98 \\
			\hline
			Variante D-FP & 10000 & 38 & 116 & 56,99 & 56 & 172,22 & 58063,98 \\
			\hline
			Variante D-CP & 100 & 31 & 72 & 47,26 & 44 & 78,93 & 1267,00 \\
			\hline
			Variante D-CP & 100 & 31 & 70 & 45,93 & 45 & 77,82 & 1285,00 \\
			\hline
			Variante D-CP & 100 & 30 & 74 & 46,81 & 46 & 80,13 & 1248,00 \\
			\hline
			Variante D-CP & 1000 & 30 & 78 & 48,06 & 47 & 180,86 & 5529,00 \\
			\hline
			Variante D-CP & 1000 & 30 & 105 & 50,88 & 49 & 170,68 & 5859,00 \\
			\hline
			Variante D-CP & 1000 & 30 & 75 & 50,10 & 50 & 173,97 & 5748,00 \\
			\hline
			Variante D-CP & 10000 & 34 & 94 & 48,87 & 49 & 199,60 & 50099,00 \\
			\hline
			Variante D-CP & 10000 & 34 & 111 & 49,95 & 50 & 195,57 & 51132,01 \\
			\hline
			Variante D-CP & 10000 & 34 & 93 & 48,28 & 48 & 200,92 & 49771,99 \\
			\hline
		\end{tabular}
		\caption{Analyseergebnis des Lasttests der verschiedenen Variationen in Kombination mit Postgres}
		\label{fig:performance-postgres}
	\end{table}
\end{landscape}

\begin{landscape}
	\begin{table}[h!]
		\centering
		\small
		\vspace{1cm}
		\begin{tabular}{ |c|c|c|c|c|c|c|c|} 
			\hline
			Name & Anzahl Durchläufe & min. AZ & max. AZ & durchschn. AZ & Median der AZ & Durchläufe pro Sekunde & Testdauer in Minuten\\ 
			\hline
			Variante A-M & 10000 & 563 & 1526 & 612,04 & 612 & 16,30 & 10,22 \\
			\hline
			Variante A-M & 10000 & 565 & 725 & 617,67 & 617 & 16,06 & 10,38 \\
			\hline
			Variante A-M & 10000 & 576 & 743 & 619,05 & 617 & 16,04 & 10,39 \\
			\hline
			Variante D-FM & 10000 & 1405 & 2904 & 1625,89 & 1612 & 6,13 & 27,19 \\
			\hline
			Variante D-FM & 10000 & 1510 & 1720 & 1608,00 & 1608 & 6,19 & 26,91 \\
			\hline
			Variante D-FM & 10000 & 1512 & 1970 & 1654,91 & 1646 & 6,01 & 27,72 \\
			\hline
			Variante D-CM & 10000 & 1210 & 3307 & 1380,43 & 1364 & 7,23 & 23,06 \\
			\hline
			Variante D-CM & 10000 & 1207 & 1675 & 1384,97 & 1368 & 7,17 & 23,26 \\
			\hline
			Variante D-CM & 10000 & 1219 & 1637 & 1374,96 & 1365 & 7,21 & 23,11 \\
			\hline
			Variante A-P & 10000 & 3354 & 4237 & 3848,23 & 3836 & 2,49 & 66,96 \\
			\hline
			Variante A-P & 10000 & 3402 & 4305 & 3769,70 & 3765 & 2,59 & 64,46 \\
			\hline
			Variante A-P & 10000 & 3364 & 4911 & 3833,21 & 3796 & 2,36 & 70,67 \\
			\hline
			Variante D-FP & 10000 & 5386 & 6919 & 5839,03 & 5804 & 1,66 & 100,49 \\
			\hline
			Variante D-FP & 10000 & 5168 & 6883 & 5778,68 & 5725 & 1,52 & 109,65 \\
			\hline
			Variante D-FP & 10000 & 5235 & 6174 & 5727,97 & 5720 & 1,67 & 99,97 \\
			\hline
			Variante D-CP & 10000 & 4189 & 5799 & 4571,81 & 4549 & 2,11 & 79,14 \\
			\hline
			Variante D-CP & 10000 & 3919 & 5553 & 4553,76 & 4554 & 2,12 & 78,61 \\
			\hline
			Variante D-CP & 10000 & 4000 & 5702 & 4544,34 & 4527 & 2,02 & 82,69 \\
			\hline
		\end{tabular}
		\caption{Lasttest-Ergebnisse mit Datenbanken von einem externen Cloud-Anbieter}
		\label{fig:performance-database}
	\end{table}
\end{landscape}



\begin{landscape}
	\begin{table}[h!]
		\centering
		\small
		\vspace{1cm}
		\begin{tabular}{ |c|c|c|c|c|c|c|c|} 
			\hline
			Name & Anzahl Durchläufe & min. AZ & max. AZ & durchschn. AZ & Median der AZ & Durchläufe pro Sekunde & Testdauer in Minuten\\ 
			\hline
			Variante A-M & 10000 & 488 & 605 & 504,39 & 503 & 19,68 & 8,47 \\
			\hline
			Variante A-M & 10000 & 489 & 585 & 502,43 & 501 & 19,71 & 8,45 \\
			\hline
			Variante A-M & 10000 & 490 & 604 & 502,69 & 501 & 19,73 & 8,45 \\
			\hline
			Variante D-FM & 10000 & 494 & 604 & 511,01 & 511 & 19,39 & 8,59 \\
			\hline
			Variante D-FM & 10000 & 494 & 600 & 509,20 & 508 & 19,38 & 8,60 \\
			\hline
			Variante D-FM & 10000 & 492 & 601 & 508,31 & 507 & 19,50 & 8,55 \\
			\hline
			Variante D-CM & 10000 & 490 & 575 & 503,62 & 502 & 19,69 & 8,46 \\
			\hline
			Variante D-CM & 10000 & 491 & 627 & 503,50 & 502 & 19,68 & 8,47 \\
			\hline
			Variante D-CM & 10000 & 491 & 571 & 505,90 & 504 & 19,61 & 8,50 \\
			\hline
			Variante A-P & 10000 & 507 & 589 & 522,61 & 522 & 18,97 & 8,78 \\
			\hline
			Variante A-P & 10000 & 507 & 630 & 519,33 & 518 & 19,08 & 8,73 \\
			\hline
			Variante A-P & 10000 & 507 & 602 & 521,18 & 520 & 19,03 & 8,76 \\
			\hline
			Variante D-FP & 10000 & 511 & 594 & 526,00 & 525 & 18,86 & 8,84 \\
			\hline
			Variante D-FP & 10000 & 511 & 600 & 526,29 & 525 & 18,82 & 8,85 \\
			\hline
			Variante D-FP & 10000 & 512 & 626 & 526,81 & 525 & 18,80 & 8,87 \\
			\hline
			Variante D-CP & 10000 & 504 & 580 & 517,00 & 516 & 19,14 & 8,71 \\
			\hline
			Variante D-CP & 10000 & 502 & 611 & 518,79 & 518 & 19,11 & 8,72 \\
			\hline
			Variante D-CP & 10000 & 505 & 594 & 518,93 & 518 & 19,09 & 8,73 \\
			\hline
		\end{tabular}
	\caption{Lasttest-Ergebnisse mit Simulation der API-Aufruf durch künstliche Verzögerung}
	\label{fig:performance-delay}
	\end{table}
\end{landscape}

\end{anhang}

  
  %Refs
  %TODO: Add alot of references
  %TODO: Add Effective Aggregate Design
  \addcontentsline{toc}{chapter}{Literatur}
  \printbibliography %Citavi 5
  
\end{document}